% !TEX Program = xelatex
\documentclass[QingNianC, vlined, ]{nsfc}



\begin{document}
	
	%%%%%%%%%%                                  %%%%%%%%%%
	%%               以下标题对应青年A类                 %%
	%%%%%%%%%%                                  %%%%%%%%%%
	\nsfctitle

	\nsfcsection{主要学术成绩、创新点及其科学意义}{（建议不超过5000字）}
	\nsfcsubsection*{}{着重阐述近5年来在基础研究方面所取得的学术成绩、创新点、研究价值和科学意义等。}

	\nsfcsection{拟开展的研究工作}{（建议不超过3000字）}
	\nsfcsubsection*{}{着重阐述拟开展的研究工作的创新性构思，主要研究方向和初步研究方案等，请简要阐述。}

	\nsfcsection{其他需要说明的情况}{}
	\nsfcsubsection{}{申请人同年申请不同类型的国家自然科学基金项目情况（列明同年申请的其他项目的项目类型、项目名称信息，并说明与本项目之间的区别与联系；已收到自然科学基金委不予受理或不予资助决定的，无需列出）。}

	\nsfcsubsection{}{具有高级专业技术职务（职称）的申请人是否存在同年申请或者参与申请国家自然科学基金项目的单位不一致的情况；如存在上述情况，列明所涉及人员的姓名，申请或参与申请的其他项目的项目类型、项目名称、单位名称、上述人员在该项目中是申请人还是参与者，并说明单位不一致原因。}

	\nsfcsubsection{}{具有高级专业技术职务（职称）的申请人是否存在与正在承担的国家自然科学基金项目的单位不一致的情况；如存在上述情况，列明所涉及人员的姓名，正在承担项目的批准号、项目类型、项目名称、单位名称、起止年月，并说明单位不一致原因。}

	\nsfcsubsection{}{同年以不同专业技术职务（职称）申请或参与申请科学基金项目的情况（应详细说明原因）。}

	\nsfcsubsection{}{其他。}
 


	%%%%%%%%%%                                  %%%%%%%%%%
	%%               以下标题对应青年B类                 %%
	%%%%%%%%%%                                  %%%%%%%%%%

	\nsfctitle
	\nsfcsection{主要学术成绩、创新点及其科学意义}{（建议不超过4000字）}
	\nsfcsubsection*{}{按本年度《国家自然科学基金项目指南》中青年科学基金项目（B类）的有关要求，着重阐述所取得的研究成果的创新性和科学价值等。}

	\nsfcsection{拟开展的研究工作}{（建议不超过4000字）}
	\nsfcsubsection*{}{着重阐述拟开展的研究工作的科学意义和创新性，技术路线、研究方案等的可行性。}

	\nsfcsection{其他需要说明的情况}{}
	\nsfcsubsection{}{申请人同年申请不同类型的国家自然科学基金项目情况（列明同年申请的其他项目的项目类型、项目名称信息，并说明与本项目之间的区别与联系；已收到自然科学基金委不予受理或不予资助决定的，无需列出）。}

	\nsfcsubsection{}{具有高级专业技术职务（职称）的申请人是否存在同年申请或者参与申请国家自然科学基金项目的单位不一致的情况；如存在上述情况，列明所涉及人员的姓名，申请或参与申请的其他项目的项目类型、项目名称、单位名称、上述人员在该项目中是申请人还是参与者，并说明单位不一致原因。}

	\nsfcsubsection{}{具有高级专业技术职务（职称）的申请人是否存在与正在承担的国家自然科学基金项目的单位不一致的情况；如存在上述情况，列明所涉及人员的姓名，正在承担项目的批准号、项目类型、项目名称、单位名称、起止年月，并说明单位不一致原因。}

	\nsfcsubsection{}{同年以不同专业技术职务（职称）申请或参与申请科学基金项目的情况（应详细说明原因）。}

	\nsfcsubsection{}{其他。}



	%%%%%%%%%%                                  %%%%%%%%%%
	%%               以下标题对应青年C类                 %%
	%%%%%%%%%%                                  %%%%%%%%%%
	\nsfctitle

	\nsfcsection{立项依据}{}

	\nsfcsubsection*{}{（为什么要开展此项研究，研究的科学技术价值如何）}

	\nsfcsection{研究内容}{}

	\nsfcsubsection*{}{（提纲不做限制，请按照研究工作的自身逻辑撰写。应提炼出特色与创新点、年度研究计划）}

	\nsfcsection{研究基础}{}

	\nsfcsubsection{研究基础与可行性分析}{（与本项目相关的研究工作积累和已取得的研究工作成绩，研究风险的应对措施等）；}

	\nsfcsubsection{工作条件}{（包括已具备的实验条件，尚缺少的实验条件和拟解决的途径，包括利用国家实验室、全国重点实验室和部门重点实验室等研究基地的计划与落实情况）；}

	\nsfcsubsection{正在承担的与本项目相关的科研项目情况}{（申请人正在承担的与本项目相关的科研项目情况，包括国家自然科学基金的项目和国家其他科技计划项目，要注明项目的资助机构、项目类别、批准号、项目名称、获资助金额、起止年月、与本项目的关系及负责的内容等）；}

	\nsfcsubsection{完成国家自然科学基金项目情况}{（对申请人负责的前一个已资助期满的科学基金项目（项目名称及批准号）完成情况、后续研究进展及与本申请项目的关系加以详细说明。另附该项目的研究工作总结摘要（限500字）和相关成果详细目录）。}

	\nsfcsection{其他需要说明的情况}{}
	\nsfcsubsection{}{申请人同年申请不同类型的国家自然科学基金项目情况（列明同年申请的其他项目的项目类型、项目名称信息，并说明与本项目之间的区别与联系；已收到自然科学基金委不予受理或不予资助决定的，无需列出）。}


	\nsfcsubsection{}{具有高级专业技术职务（职称）的申请人是否存在同年申请或者参与申请国家自然科学基金项目的单位不一致的情况；如存在上述情况，列明所涉及人员的姓名，申请或参与申请的其他项目的项目类型、项目名称、单位名称、上述人员在该项目中是申请人还是参与者，并说明单位不一致原因。}

	\nsfcsubsection{}{具有高级专业技术职务（职称）的申请人是否存在与正在承担的国家自然科学基金项目的单位不一致的情况；如存在上述情况，列明所涉及人员的姓名，正在承担项目的批准号、项目类型、项目名称、单位名称、起止年月，并说明单位不一致原因。}

	\nsfcsubsection{}{同年以不同专业技术职务（职称）申请或参与申请科学基金项目的情况（应详细说明原因）。}

	\nsfcsubsection{}{其他。}



	%%%%%%%%%%                                  %%%%%%%%%%
	%%               以下标题对应面上项目                 %%
	%%%%%%%%%%                                  %%%%%%%%%%

	\nsfctitle
	
	\nsfcsection{立项依据：}{}
	\nsfcsubsection*{}{（为什么要开展此项研究，研究的科学技术价值如何）}

	\nsfcsection{研究内容：}{}
	\nsfcsubsection*{}{（提纲不做限制，请按照研究工作的自身逻辑撰写。应提炼出特色与创新点、年度研究计划）}

	\nsfcsection{研究基础：}{}
	\nsfcsubsection{研究基础与可行性分析}{（与本项目相关的研究工作积累和已取得的研究工作成绩，研究风险的应对措施等）；}

	\nsfcsubsection{工作条件}{（包括已具备的实验条件，尚缺少的实验条件和拟解决的途径，包括利用国家实验室、全国重点实验室和部门重点实验室等研究基地的计划与落实情况）；}

	\nsfcsubsection{正在承担的与本项目相关的科研项目情况}{（申请人和主要参与者正在承担的与本项目相关的科研项目情况，包括国家自然科学基金的项目和国家其他科技计划项目，要注明项目的资助机构、项目类别、批准号、项目名称、获资助金额、起止年月、与本项目的关系及负责的内容等）；}

	\nsfcsubsection{完成国家自然科学基金项目情况}{（对申请人负责的前一个已资助期满的科学基金项目（项目名称及批准号）完成情况、后续研究进展及与本申请项目的关系加以详细说明。另附该项目的研究工作总结摘要（限500字）和相关成果详细目录）。}

	\nsfcsection{其他需要说明的情况：}{}
	\nsfcsubsection{}{申请人同年申请不同类型的国家自然科学基金项目情况（列明同年申请的其他项目的项目类型、项目名称信息，并说明与本项目之间的区别与联系；已收到自然科学基金委不予受理或不予资助决定的，无需列出）。}

	\nsfcsubsection{}{具有高级专业技术职务（职称）的申请人或者主要参与者是否存在同年申请或者参与申请国家自然科学基金项目的单位不一致的情况；如存在上述情况，列明所涉及人员的姓名，申请或参与申请的其他项目的项目类型、项目名称、单位名称、上述人员在该项目中是申请人还是参与者，并说明单位不一致原因。}

	\nsfcsubsection{}{具有高级专业技术职务（职称）的申请人或者主要参与者是否存在与正在承担的国家自然科学基金项目的单位不一致的情况；如存在上述情况，列明所涉及人员的姓名，正在承担项目的批准号、项目类型、项目名称、单位名称、起止年月，并说明单位不一致原因。}

	\nsfcsubsection{}{申请人和主要参与者同年以不同专业技术职务（职称）申请或参与申请科学基金项目的情况（应详细说明原因）。}

	\nsfcsubsection{}{其他。}
 



	%%%%%%%%%%                                  %%%%%%%%%%
	%%               以下标题对应重点项目                 %%
	%%%%%%%%%%                                  %%%%%%%%%%

	\nsfctitle

	\nsfcsection{立项依据与研究内容（建议8000字以下）}{}
	\nsfcsubsection{项目的立项依据}{（研究意义、国内外研究现状及发展动态分析，需结合科学研究发展趋势来论述科学意义；或结合国民经济和社会发展中迫切需要解决的关键科技问题来论述其应用前景。附主要参考文献目录）；}

	\nsfcsubsection{项目的研究内容、研究目标，以及拟解决的关键科学问题}{（此部分为重点阐述内容）；}

	\nsfcsubsection{拟采取的研究方案及可行性分析}{（包括研究方法、技术路线、实验手段、关键技术等说明）；}

	\nsfcsubsection{本项目的特色与创新之处；}{}

	\nsfcsubsection{年度研究计划及预期研究结果}{（包括拟组织的重要学术交流活动、国际合作与交流计划等）。}

	\nsfcsection{研究基础与工作条件}{}
	\nsfcsubsection{研究基础}{（与本项目相关的研究工作积累和已取得的研究工作成绩）；}

	\nsfcsubsection{工作条件}{（包括已具备的实验条件，尚缺少的实验条件和拟解决的途径，包括利用国家实验室、全国重点实验室和部门重点实验室等研究基地的计划与落实情况）；}

	\nsfcsubsection{正在承担的与本项目相关的科研项目情况}{（申请人和主要参与者正在承担的与本项目相关的科研项目情况，包括国家自然科学基金的项目和国家其他科技计划项目，要注明项目的资助机构、项目类别、批准号、项目名称、获资助金额、起止年月、与本项目的关系及负责的内容等）；}

	\nsfcsubsection{完成国家自然科学基金项目情况}{（对申请人负责的前一个已资助期满的科学基金项目（项目名称及批准号）完成情况、后续研究进展及与本申请项目的关系加以详细说明。另附该项目的研究工作总结摘要（限500字）和相关成果详细目录）。}

	\nsfcsection{其他需要说明的情况}{}
	\nsfcsubsection{}{申请人同年申请不同类型的国家自然科学基金项目情况（列明同年申请的其他项目的项目类型、项目名称信息，并说明与本项目之间的区别与联系；已收到自然科学基金委不予受理或不予资助决定的，无需列出）。}

	\nsfcsubsection{}{具有高级专业技术职务（职称）的申请人或者主要参与者是否存在同年申请或者参与申请国家自然科学基金项目的单位不一致的情况；如存在上述情况，列明所涉及人员的姓名，申请或参与申请的其他项目的项目类型、项目名称、单位名称、上述人员在该项目中是申请人还是参与者，并说明单位不一致原因。}

	\nsfcsubsection{}{具有高级专业技术职务（职称）的申请人或者主要参与者是否存在与正在承担的国家自然科学基金项目的单位不一致的情况；如存在上述情况，列明所涉及人员的姓名，正在承担项目的批准号、项目类型、项目名称、单位名称、起止年月，并说明单位不一致原因。}

	\nsfcsubsection{}{申请人和主要参与者同年以不同专业技术职务（职称）申请或参与申请科学基金项目的情况（应详细说明原因）。}

	\nsfcsubsection{}{其他。}




	%%%%%%%%%%                                  %%%%%%%%%%
	%%               以下标题对应重大项目                 %%
	%%%%%%%%%%                                  %%%%%%%%%%

	\nsfctitle

	\nsfcsection{立项依据与研究内容}{（建议在5000-10000字之间）：}
	\nsfcsubsection{项目的立项依据}{（研究意义、国内外研究现状及发展动态分析，需结合科学研究发展趋势来论述科学意义；或结合国民经济和社会发展中迫切需要解决的关键科技问题来论述其应用前景。附主要参考文献目录）；}

	\nsfcsubsection{项目的研究内容、研究目标以及拟解决的关键科学问题}{（此部分为重点阐述内容）；}

	\nsfcsubsection{总体研究方案、学术思路、技术途径及可行性分析；}{}

	\nsfcsubsection{课题设置}{（所设各课题的研究重点和彼此分工，达到预期目标需安排的交叉、综合研究和联合实验、观测研究）；}

	\nsfcsubsection{年度研究计划}{（包括拟组织的重要学术交流活动、国际合作与交流计划等）；}

	\nsfcsubsection{预期研究成果及可能达到的水平}{（理论性研究成果需注明预期水平和影响；应用性研究成果需注明其应用前景、能达到的技术指标、潜在的经济与社会效益等）。}

	\nsfcsection{项目下设课题简介：}{}
	\nsfcsubsection*{}{主要内容包括：课题名称、主要研究内容和预期目标、申请人、依托单位、合作研究单位。}

	\nsfcsection{研究基础与工作条件：}{}
	\nsfcsubsection{研究基础}{（与本项目相关的研究工作积累和已取得的研究工作成绩）；}

	\nsfcsubsection{工作条件}{（包括已具备的实验条件、基础数据资料，尚缺少的实验条件和拟解决的途径，包括利用国家实验室、国家重点实验室和部门重点实验室等研究基地的计划与落实情况）；}

	\nsfcsubsection{正在承担的与本项目相关的科研项目情况}{（申请人和主要参与者正在承担的与本项目相关的科研项目情况，包括国家自然科学基金的项目和国家其他科技计划项目，要注明项目的资助机构、项目类别、批准号、项目名称、获资助金额、起止年月、与本项目的关系及负责的内容等）；}

	\nsfcsubsection{完成国家自然科学基金项目情况}{（对申请人负责的前一个已资助期满的科学基金项目（项目名称及批准号）完成情况、后续研究进展及与本申请项目的关系加以详细说明。另附该项目的研究工作总结摘要（限500字）和相关成果详细目录）。}

	\nsfcsection{其他需要说明的情况：}{}
	\nsfcsubsection{}{申请人同年申请不同类型的国家自然科学基金项目情况（列明同年申请的其他项目的项目类型、项目名称信息，并说明与本项目之间的区别与联系;已收到自然科学基金委不予受理或不予资助决定的，无需列出）。}

	\nsfcsubsection{}{具有高级专业技术职务（职称）的申请人或者主要参与者是否存在同年申请或者参与申请国家自然科学基金项目的单位不一致的情况；如存在上述情况，列明所涉及人员的姓名，申请或参与申请的其他项目的项目类型、项目名称、单位名称、上述人员在该项目中是申请人还是参与者，并说明单位不一致原因。}

	\nsfcsubsection{}{具有高级专业技术职务（职称）的申请人或者主要参与者是否存在与正在承担的国家自然科学基金项目的单位不一致的情况；如存在上述情况，列明所涉及人员的姓名，正在承担项目的批准号、项目类型、项目名称、单位名称、起止年月，并说明单位不一致原因。}

	\nsfcsubsection{}{申请人和主要参与者同年以不同专业技术职务（职称）申请或参与申请科学基金项目的情况（应详细说明原因）。}

	\nsfcsubsection{}{其他。}




	%%%%%%%%%%                                  %%%%%%%%%%
	%%               以下标题对应重大课题                 %%
	%%%%%%%%%%                                  %%%%%%%%%%

	\nsfctitle

	\nsfcsection{立项依据与研究内容}{（建议在5000-10000字之间）：}
	\nsfcsubsection{课题的立项依据}{（研究意义、国内外研究现状及发展动态分析，需结合科学研究发展趋势来论述科学意义；或结合国民经济和社会发展中迫切需要解决的关键科技问题来论述其应用前景。附主要参考文献目录）；}

	\nsfcsubsection{课题的研究内容、研究目标、拟解决的关键科学问题及在重大项目中所起的作用}{（此部分为重点阐述内容）；}

	\nsfcsubsection{拟采取的研究方案及可行性分析}{（包括研究方法、技术路线、实验手段、关键技术等说明）；}

	\nsfcsubsection{课题的特色与创新之处；}

	\nsfcsubsection{年度研究计划及预期研究结果}{（包括拟组织的重要学术交流活动、国际合作与交流计划等）。}

	\nsfcsection{研究基础与工作条件：}{}
	\nsfcsubsection{研究基础}{（与本课题相关的研究工作积累和已取得的研究工作成绩）；}

	\nsfcsubsection{工作条件}{（包括已具备的实验条件，尚缺少的实验条件和拟解决的途径，包括利用国家实验室、国家重点实验室和部门重点实验室等研究基地的计划与落实情况）；}

	\nsfcsubsection{正在承担的与本项目相关的科研项目情况}{（申请人和主要参与者正在承担的与本课题相关的科研项目情况，包括国家自然科学基金的项目和国家其他科技计划项目，要注明项目的资助机构、项目类别、批准号、项目名称、获资助金额、起止年月、与本项目的关系及负责的内容等）；}

	\nsfcsubsection{完成国家自然科学基金项目情况}{（对申请人负责的前一个已资助期满的科学基金项目（项目名称及批准号）完成情况、后续研究进展及与本申请项目的关系加以详细说明。另附该项目的研究工作总结摘要（限500字）和相关成果详细目录）。}

	\nsfcsection{其他需要说明的情况：}{}
	\nsfcsubsection{}{申请人同年申请不同类型的国家自然科学基金项目情况（列明同年申请的其他项目的项目类型、项目名称信息，并说明与本项目之间的区别与联系;已收到自然科学基金委不予受理或不予资助决定的，无需列出）。}

	\nsfcsubsection{}{具有高级专业技术职务（职称）的申请人或者主要参与者是否存在同年申请或者参与申请国家自然科学基金项目的单位不一致的情况；如存在上述情况，列明所涉及人员的姓名，申请或参与申请的其他项目的项目类型、项目名称、单位名称、上述人员在该项目中是申请人还是参与者，并说明单位不一致原因。}

	\nsfcsubsection{}{具有高级专业技术职务（职称）的申请人或者主要参与者是否存在与正在承担的国家自然科学基金项目的单位不一致的情况；如存在上述情况，列明所涉及人员的姓名，正在承担项目的批准号、项目类型、项目名称、单位名称、起止年月，并说明单位不一致原因。}

	\nsfcsubsection{}{申请人和主要参与者同年以不同专业技术职务（职称）申请或参与申请科学基金项目的情况（应详细说明原因）。}

	\nsfcsubsection{}{其他。}




	%%%%%%%%%%                                  %%%%%%%%%%
	%%               以下标题对应重大计划                 %%
	%%%%%%%%%%                                  %%%%%%%%%%

	\nsfctitle

	\nsfcsection{立项依据与研究内容}{（建议在5000-10000字之间）：}
	\nsfcsubsection{项目的立项依据}{（研究意义、国内外研究现状及发展动态分析，需结合科学研究发展趋势来论述科学意义；或结合国民经济和社会发展中迫切需要解决的关键科技问题来论述其应用前景。附主要参考文献目录）；}

	\nsfcsubsection{项目的研究内容、研究目标，以及拟解决的关键科学问题}{（此部分为重点阐述内容；另外，请说明本项目与重大研究计划总体目标的关系）；}

	\nsfcsubsection{拟采取的研究方案及可行性分析}{（包括研究方法、技术路线、实验手段、关键技术等说明）；}

	\nsfcsubsection{本项目的特色与创新之处；}{}

	\nsfcsubsection{年度研究计划及预期研究结果}{（包括拟组织的重要学术交流活动、国际合作与交流计划等）。}

	\nsfcsection{研究基础与工作条件：}{}
	\nsfcsubsection{研究基础}{（与本项目相关的研究工作积累和已取得的研究工作成绩）；}

	\nsfcsubsection{工作条件}{（包括已具备的实验条件，尚缺少的实验条件和拟解决的途径，包括利用国家实验室、国家重点实验室和部门重点实验室等研究基地的计划与落实情况）；}

	\nsfcsubsection{正在承担的与本项目相关的科研项目情况}{（申请人和主要参与者正在承担的与本项目相关的科研项目情况，包括国家自然科学基金的项目和国家其他科技计划项目，要注明项目的资助机构、项目类别、批准号、项目名称、获资助金额、起止年月、与本项目的关系及负责的内容等）；}

	\nsfcsubsection{完成国家自然科学基金项目情况}{（对申请人负责的前一个已资助期满的科学基金项目（项目名称及批准号）完成情况、后续研究进展及与本申请项目的关系加以详细说明。另附该项目的研究工作总结摘要（限500字）和相关成果详细目录）。}

	\nsfcsection{其他需要说明的情况：}{}
	\nsfcsubsection{}{申请人同年申请不同类型的国家自然科学基金项目情况（列明同年申请的其他项目的项目类型、项目名称信息，并说明与本项目之间的区别与联系；已收到自然科学基金委不予受理或不予资助决定的，无需列出）。}

	\nsfcsubsection{}{具有高级专业技术职务（职称）的申请人或者主要参与者是否存在同年申请或者参与申请国家自然科学基金项目的单位不一致的情况；如存在上述情况，列明所涉及人员的姓名，申请或参与申请的其他项目的项目类型、项目名称、单位名称、上述人员在该项目中是申请人还是参与者，并说明单位不一致原因。}

	\nsfcsubsection{}{具有高级专业技术职务（职称）的申请人或者主要参与者是否存在与正在承担的国家自然科学基金项目的单位不一致的情况；如存在上述情况，列明所涉及人员的姓名，正在承担项目的批准号、项目类型、项目名称、单位名称、起止年月，并说明单位不一致原因。}

	\nsfcsubsection{}{申请人和主要参与者同年以不同专业技术职务（职称）申请或参与申请科学基金项目的情况（应详细说明原因）。}

	\nsfcsubsection{}{其他。}



	%%%%%%%%%%%%%%%%%%%%%%%%%%%%%%%%%%%%%%%%%%%%%%%%%%%%%%%%%%%%%%%
	%%     以下为模版说明及使用示例，挑需要的用，定稿时删掉或注释     %%
	%%%%%%%%%%%%%%%%%%%%%%%%%%%%%%%%%%%%%%%%%%%%%%%%%%%%%%%%%%%%%%%

	\newpage
	\stepcounter{chapter}

	% 以下两条命令为高亮示例文档中的关键内容而设，正式撰写时不要使用
	\newcommand{\shad}[1]{\textcolor{DodgerBlue}{\ttfamily #1}}
	\newcommand{\shadcmd}[1]{\shad{$\backslash$#1}}


	\section{前言}
	\subsection{当前版本编译时间}
	\noindent%
	\begin{center}%
		\zihao{-2}\bfseries\color{DarkRed}%
		\zhdigits*{\the\year} 年\zhnumber{\the\month} 月\zhnumber{\the\day} 日\;\xxivtime
	\end{center}

	\subsection{适用对象}
	
	本模板参照自然科学基金委官方提供的Word模版编写，适用于\textbf{青年基金A}、\textbf{青年基金B}、\textbf{青年基金C}、\textbf{面上项目}、\textbf{重点项目}、\textbf{重大项目}、\textbf{重大项目课题}和\textbf{重大研究计划}。

	\subsection{使用环境}

	本模板兼容Windows、MacOS以及Overleaf等主流平台，其开发测试环境为TeXLive2026+VSCode。\textbf{建议本地用户将LaTeX环境更新到\href{https://mirrors.tuna.tsinghua.edu.cn/CTAN/systems/texlive/Images/}{\color{DarkRed}\CJKunderline*[thickness=0.5bp, format=\color{DarkRed}]{TeXLive2026及以上}}或\href{https://mirrors.tuna.tsinghua.edu.cn/tex-historic-archive/systems/mactex/2026/}{\color{DarkRed}\CJKunderline*[thickness=0.5bp, format=\color{DarkRed}]{MacTeX2026及以上}}}（至少不低于\href{https://mirrors.tuna.tsinghua.edu.cn/tex-historic-archive/systems/texlive/}{\color{DarkRed}\CJKunderline*[thickness=0.5bp, format=\color{DarkRed}]{TeXLive2024}}或\href{https://mirrors.tuna.tsinghua.edu.cn/tex-historic-archive/systems/mactex/}{\color{DarkRed}\CJKunderline*[thickness=0.5bp, format=\color{DarkRed}]{MacTeX2024}}），以避免兼容性问题。

	\begin{itemize}
		\item \href{https://zhuanlan.zhihu.com/p/389394015}{\color{DarkRed}\CJKunderline*[thickness=0.5bp, format=\color{DarkRed}]{TeXLive安装教程}}、\href{https://blog.csdn.net/ChrisP_333/article/details/82943508}{\color{DarkRed}\CJKunderline*[thickness=0.5bp, format=\color{DarkRed}]{MacTeX安装教程}}
		\item \href{https://code.visualstudio.com/Download}{\color{DarkRed}\CJKunderline*[thickness=0.5bp, format=\color{DarkRed}]{VSCode下载地址}}、\href{https://zhuanlan.zhihu.com/p/166523064}{\color{DarkRed}\CJKunderline*[thickness=0.5bp, format=\color{DarkRed}]{VSCode配置LaTeX环境}}、\href{https://texstudio.sourceforge.net/#download}{\color{DarkRed}\CJKunderline*[thickness=0.5bp, format=\color{DarkRed}]{TeXstudio下载地址}}
	\end{itemize}
	
	模板需要使用\shad{XeLaTeX}引擎编译：
	\begin{itemize}
		\item 若使用TeXstudio编辑器，则不需要对软件进行设置：\shad{tutorial.tex}文件首行已经添加了\shad{\% !TEX Program = xelatex}，该指令指定使用\shad{XeLaTeX}编译该文档；
		
		\item VSCode编辑器和Overleaf需要自行将编译引擎设置为\shad{XeLaTeX}。
	\end{itemize}
	

	为了确保在Windows、MacOS、Overleaf等平台上编译出完全一致的结果，模板内置了所有用到的字体文件，项目大小可能超出Overleaf上传压缩包的限制。若遇到这种情况，\textbf{先在Overleaf上新建一个空项目，然后解压本模板并拖拽文件和文件夹到新建的项目中}。\textcolor{DarkRed}{\textbf{注意在网站里切换到TeXLive2026}}。

	\subsection{模板更新方法}

	模板的发布地址为：\href{https://github.com/MGG1996/NSFC}{\color{DarkRed}\CJKunderline*[thickness=0.5bp, format=\color{DarkRed}]{https://github.com/MGG1996/NSFC}}。
	
	更新模板的正确方式：\textbf{下载最新的完整压缩包，解压后用自己的\shad{.bib}和\shad{.tex}文件以及\shad{fig}目录替换掉模板中原有的同名文件和目录}。

	更建议的更新方式：\textbf{\color{DarkRed}在初次使用本模板时，修改\shad{.tex}和\shad{.bib}文件的文件名（别忘了\shad{.tex}文件内通过\shadcmd{bibliography\{\}}设置的\shad{.bib}文件名），后续只需下载最新的模板，并将其内容全部复制到申报书所在目录进行替换}。


	\subsection{新手入门}

	使用本模板需要掌握基本的LaTeX排版操作。如果你是纯新手，那可以先看看\href{https://ctan.math.utah.edu/ctan/tex-archive/info/lshort/chinese/lshort-zh-cn.pdf}{\bfseries\color{DarkRed}\CJKunderline*[thickness=0.5bp, format=\color{DarkRed}]{一份（不太）简短的LaTeX2ε介绍}}，它足以帮助你建立起基本概念。


	\subsection{本模板已载入的宏包}

	本模板已在\shad{nsfc.cls}文件中载入下列宏包，用户应避免重复载入，有特殊需求的用户在载入其他宏包时应留意与现有宏包是否冲突。

	\begin{verbatim}
		\RequirePackage{ifthen}  % 条件判断需要
		\RequirePackage[UTF8]{ctex}  % ctex 包含 xeCJK 包含 fontspec
		\RequirePackage{xeCJKfntef}  % 中文添加下划线需要
		\RequirePackage{calc}  % 简单的长度加减运算需要
		\RequirePackage{color}
		\RequirePackage[dvipsnames, svgnames, x11names]{xcolor}  % 更多预设颜色
		\RequirePackage{layouts}  % 设置页面布局
		\RequirePackage{zhnumber}  % 将计数器显示成中文
		\RequirePackage{titlesec}  % 修改各级标题样式、目录等
		\RequirePackage{titletoc}  % 调整目录需要
		\RequirePackage[titles]{tocloft}  % 调整图目录、表目录需要
		\RequirePackage[
		    bookmarksnumbered=true,
		    bookmarksdepth=3,
		    citecolor=black,
		    linkcolor=black,
		    urlcolor=black,
		]{hyperref}  % 生成带编号的书签，并设置书签深度
		\RequirePackage{graphicx}  % 图片排版
		\RequirePackage[font=small]{subfig}  % 子图排版
		\RequirePackage[numbers,sort&compress]{natbib}  % 参考文献及其引用样式
		\RequirePackage{amssymb, amsmath, amsthm}  % 公式、数学符号、定理环境
		\RequirePackage{caption}  % 图、表、伪代码标题
		\RequirePackage[algochapter, linesnumbered, ruled]{algorithm2e}% 伪码
		\RequirePackage{xunicode-addon}  % 脚注编号使用带圈数字需要
		\RequirePackage{xpatch}
		\RequirePackage{scrextend}  % 调整脚注的缩进设置
		\RequirePackage[perpage]{footmisc}  % 脚注序号每页重置
		\RequirePackage{enumitem}  % 调整list的边距需要
		\RequirePackage{array}
		\RequirePackage[flushleft]{threeparttable}  % 排版带附注的表格需要
		\RequirePackage{booktabs}  % 调整三线表线宽需要
		\RequirePackage{longtable}  % 跨页表格需要
		\RequirePackage{xparse}  % 定义带多个可选参数的命令或环境时需要
		\RequirePackage{multirow}  % 表格合并多行单元格需要
		\RequirePackage{extarrows}  % 带文字的长等号与箭头需要
		\RequirePackage{mathspec}  % 设置公式字体为Times New Roman需要
		\RequirePackage{xstring}  % 在字符串中搜索子串需要
		\RequirePackage{listings}  % 代码块风格化排版需要
		\RequirePackage{mhchem, chemfig}  % 排版化学方程式、结构式和键线式
		\RequirePackage{datetime}  % 获取当前系统时间需要
	\end{verbatim}
	

	\section{模板使用说明}
	\subsection{导言区及模板选项}

	模板的导言区只有两行：
	\begin{itemize}
		\item \shad{\% !TEX Program = xelatex}：在Texstudio中，这表示指定使用\shad{XeLaTeX}引擎编译该文档；在其他编辑器中，需要手动设置编译引擎为\shad{XeLaTeX}。

		\item \shadcmd{documentclass[<选项列表>]\{nsfc\}}：加载名为\shad{nsfc} 的文档类，该文档类基于LaTeX的\shad{book}类编写，可设置\textbf{\textcolor{DarkRed}{四种}}选项：

		\begin{itemize}
			\item \shad{QingNianA} / \shad{QingNianB} / \shad{QingNianC}\textcolor{DarkRed}{〔默认〕} / \shad{MianShang} / \shad{ZhongDian} / \shad{ZhongDa} / \shad{ZhongDaKeTi} / \shad{ZhongDaJiHua}：该选项设置基金类型，将控制页边距、正文开头的总标题等。其与基金的对应关系如下：
			
			\begin{table}[!htb]
				\centering\hspace*{1.5cm}
				\begin{tabular}{p{6em} p{9em} p{7em} p{7em}}
					\toprule
					\textbf{选项} & \textbf{对应基金类型} & \textbf{选项} & \textbf{对应基金类型}\\
					\midrule
					\shad{QingNianA} & 青年A & \shad{ZhongDian} & 重点项目 \\
					\shad{QingNianB} & 青年B & \shad{ZhongDa} & 重大项目 \\
					\shad{QingNianC} & 青年C & \shad{ZhongDaKeTi} & 重大课题 \\
					\shad{MianShang} & 面上项目 & \shad{ZhongDaJiHua} & 重大研究计划 \\
					\bottomrule
				\end{tabular}
			\end{table}
   
			\item \shad{subfigsimple} / \shad{subfigparens}\textcolor{DarkRed}{〔默认〕}：该选项用于调整正文中对子图标签进行引用时生成的编号样式，\shad{subfigsimple}对应样式\textcolor{DodgerBlue}{1-1a}；\shad{subfigparens}对应样式\textcolor{DodgerBlue}{1-1(a)}。
			
			\item \shad{cmmmath}\textcolor{DarkRed}{〔默认〕} / \shad{timesmathnogreek} / \shad{timesmath}：该选项用于选择渲染公式使用的字体。其中，\shad{cmmmath}即对应LaTeX默认使用的Computer Modern Math，也是此类选项的默认值；\shad{timesmathnogreek}指定使用Times New Roman来渲染公式中的英文字母和数字，但不影响希腊字母、手写体和双线体；\shad{timesmath}则继续将希腊字母也设置成Times New Roman（个人觉得有些违和），手写体和双线体仍保持原样。
			
			后两种选项均基于\href{https://mirrors.pku.edu.cn/ctan/macros/xetex/latex/mathspec/mathspec.pdf}{\ttfamily\color{DarkRed}\CJKunderline*[thickness=0.5bp, format=\color{DarkRed}]{mathspec}}宏包实现，经测试，只有它能做到真正意义上的Times New Roman。

			\textbf{P.s：}因为\href{https://mirrors.pku.edu.cn/ctan/macros/xetex/latex/mathspec/mathspec.pdf}{\ttfamily\color{DarkRed}\CJKunderline*[thickness=0.5bp, format=\color{DarkRed}]{mathspec}}宏包本身的特性，使用Times New Roman作为公式字体需要用户付出更多精力。例如，你想排版\shad{\$f\^{}t\$}，那么你会发现\shad{f}和\shad{t}之间的间隔很小，两者发生了重叠，这时候需要你手动用\shad{"}插入空格，变成\shad{\$f\^{}\{"t\}\$}。因此，\shad{timesmathnogreek}和\shad{timesmath}选项均存在类似瑕疵，届时就需要仔细查阅\href{https://mirrors.pku.edu.cn/ctan/macros/xetex/latex/mathspec/mathspec.pdf}{\ttfamily\color{DarkRed}\CJKunderline*[thickness=0.5bp, format=\color{DarkRed}]{mathspec}}的宏包文档。

			\item 另外，\href{https://mirrors.sustech.edu.cn/CTAN/macros/latex/contrib/algorithm2e/doc/algorithm2e.pdf}{\ttfamily\color{DarkRed}\CJKunderline*[thickness=0.5bp, format=\color{DarkRed}]{algorithm2e}}宏包的\shad{vlined}和\shad{boxruled}选项也能通过文档类设置。
		\end{itemize}
	\end{itemize}
	
	% \newpage
	\subsection{各级标题}
	
	本模板基于\shad{book}类，正文开头的总标题使用\shadcmd{nsfctitle}生成，其内容根据设置的基金类型选项自动确定。模板原有的两级\shad{蓝色标题}分别用：\par
	\shadcmd{nsfcsection\{<加粗部分>\}\{<非加粗部分>\}}\par
	\shadcmd{nsfcsubsection\{<加粗部分>\}\{<非加粗部分>\}}\\
	上述命令将自动产生编号；若无需编号，则用对应命令的带\shad{*}版本。
	
	创建你自己的各级标题则仍使用：\par
	\shad{$\backslash$subsection\{<1.型标题>\}}\par
	\shad{$\backslash$subsubsection \{<1.1型标题>\}}\par
	\shad{$\backslash$paragraph\{<（1）型标题>\}}

	\subsection{图片}
	
	本模板使用\href{https://mirror.nyist.edu.cn/CTAN/macros/latex/required/graphics/grfguide.pdf}{\ttfamily\color{DarkRed}\CJKunderline*[thickness=0.5bp, format=\color{DarkRed}]{graphicx}}和\href{https://mirrors.bfsu.edu.cn/CTAN/macros/latex/contrib/subfig/subfig.pdf}{\ttfamily\color{DarkRed}\CJKunderline*[thickness=0.5bp, format=\color{DarkRed}]{subfig}}宏包来处理插入的图片及子图，需要将待排版图片文件放入项目目录\shad{./fig/}中。以下是一些排版图片的例子。
	

	\begin{figure}[!h]
		\centering
		\includegraphics[width=0.5\linewidth]{huangrongguojing1}
		\caption{锁定目标}
	\end{figure}

	需要注意，在图\ref{fig: 报仇哪有姑姑重要}的图题中引用子图\ref{fig: 见到姑姑嘻嘻}的编号和在本段中引用子图编号使用的命令分别为\shad{$\backslash$subref\{fig: 见到姑姑嘻嘻\}}和\shad{$\backslash$ref\{fig: 见到姑姑嘻嘻\}}，它们分别生成仅含带括号子图编号和完整子图编号的结果，即\subref{fig: 见到姑姑嘻嘻}和\ref{fig: 见到姑姑嘻嘻}。
	
	\begin{figure}[!h]
		\centering
		\subfloat[]{
			\includegraphics[width=0.3\linewidth]{yangguoxiaolongnv3}
			\label{fig: 见到姑姑嘻嘻}
		}
		\hspace*{0.1\linewidth}
		\subfloat[]{
			
\includegraphics[width=0.3\linewidth]{yangguoxiaolongnv6}
			\label{fig: 姑姑见我不嘻嘻}
		}
		\caption[报仇哪有姑姑重要]{报仇哪有姑姑重要。\subref{fig: 见到姑姑嘻嘻}见到姑姑我嘻嘻；\subref{fig: 姑姑见我不嘻嘻}姑姑见我不嘻嘻} \label{fig: 报仇哪有姑姑重要}
	\end{figure}
	
	\begin{figure}[!htb]
		\centering
		\subfloat[]{
			
\includegraphics[width=0.4\linewidth]{luwushuang2}
			\label{fig: 陆无双2}
		}
		\hfill
		\subfloat[]{
			
\includegraphics[width=0.4\linewidth]{chengying3}
			\label{fig: 程英3}
		}
		\caption{找其他红颜知己嘻嘻。\subref{fig: 陆无双2}眼睛像姑姑；\subref{fig: 程英3}举止像姑姑} \label{fig: 红颜知己}
	\end{figure}
	
	此外，多行图题左右侧的缩进距离设置为单侧\shad{4em}，见图\ref{fig: 被动技能}。

	\begin{figure}[!htb]
		\centering
		\subfloat[]{
			
\includegraphics[width=0.3\linewidth]{lve2}
			\label{fig: 绿萼2}
		}
		\hspace*{0.1\linewidth}
		\subfloat[]{
			
\includegraphics[width=0.3\linewidth]{yangguolve}
			\label{fig: 杨过绿萼}
		}
		\\
		\subfloat[]{
			
\includegraphics[width=0.7\linewidth]{luwushuangchengying}
			\label{fig: GG}
		}
		\caption{撩妹是我杨过的被动技能。\subref{fig: 绿萼2}好腼腆的姑娘；\subref{fig: 杨过绿萼}你终于肯笑了；\subref{fig: GG}哦吼} \label{fig: 被动技能}
	\end{figure}
	
	\clearpage
	\subsection{表格}

	\textbf{写在最开始}：由于一些实现上的问题，对于\textbf{确定非置底排版}的任何表格，用户在通过\shad{table}环境的选项指定可选择的排版模式时，\textbf{\color{DarkRed}不可提供\shad{b}模式}，否则表格的上间距将过窄；反之，对于页内\textbf{确定置底排版}的第一个表格，用户需要\textbf{\color{DarkRed}显式在选项中指定\shad{b}或\shad{!b}}，否则此表格上间距将过宽。

	若出现意料之外的情况，用户可以通过在\shad{table}环境开始后和结束前的位置插入\shad{$\backslash$vspace*\{<距离长度>\}}来调整其上下间距，使之看起来协调。
	
	\subsubsection{普通表格}
	
	普通表格的排版本身无需多言，使用\shad{table}+\shad{tabular}环境即可，但是要注意三线表中的三条线分别需要使用\shad{$\backslash$toprule}、\shad{$\backslash$midrule}、\shad{$\backslash$bottomrule}生成，相应线粗分别为\shad{1.5}磅、\shad{0.75}磅、\shad{1.5}磅。注意不要用\shad{$\backslash$hline}，它无法调整线粗。
	
	\begin{table}[htp]
		\caption{$\backslash$cmidrule示例}\label{tab: cmidrule示例}
		\begin{threeparttable}
			\begin{tabular}{ccccc}
				\toprule
				\multirow{2}{*}{Column0} &  \multicolumn{2}{c}{Column1\tnote{1}} & \multicolumn{2}{c}{Column2\tnote{2}} \\
				\cmidrule(lr){2-3}\cmidrule[2.5bp](l){4-5}
				~     & subcolumn1 & subcolumn2 & subcolumn1 & subcolumn2 \\
				\midrule
				Row1  & element11 & element12 &element13 & element14 \\
				Row2  & element21 & element22 &element23 & element24 \\
				\cmidrule{2-3}\cmidrule[2.5bp]{4-5}
				Row3  & element31 & element32 &element33 & element34 \\
				\bottomrule
			\end{tabular}
			\begin{tablenotes}
				\item[1] \shadcmd{cmidrule}默认线粗设置为0.75bp，可用第二项可选参数同时剪去框线的左右端点
				\item[2] 通过指定\shadcmd{cmidrule}的第一项可选参数调整线粗，并用第二项可选参数仅修剪掉框线的左端点
			\end{tablenotes}
		\end{threeparttable}
	\end{table}

	在需要为表格中的某些单元格添加水平框线时，可使用\newline\shadcmd{cmidrule[<线粗>](<修剪>)\{<起始列-终止列>\}}。其第一项可选参数用于设置框线粗细，默认值设置为\shad{0.75}磅。第二项可选参数用于对框线的端点进行修剪，以防止同行独立的相邻框线在视觉上连通到一起，参见表\ref{tab: cmidrule示例}中的示例。有关第二项可选参数的取值，请查阅\href{https://mirrors.sustech.edu.cn/CTAN/macros/latex/contrib/booktabs/booktabs.pdf}{\ttfamily\color{DarkRed}\CJKunderline*[thickness=0.5bp, format=\color{DarkRed}]{booktabs}}宏包的官方文档。
	
	若要调整表格中单元格的水平对齐方式，自适应宽度的表格可以直接使用\shad{tabular}环境的\shad{l/c/r}参数；如果同时需要限定单元格宽度，则可使用\shad{p\{<宽度值>\}}、\shad{p\{<宽度值>\}<\{$\backslash$centering\}}、\shad{p\{<宽度值>\}<\{$\backslash$raggedleft\}}参数，效果见表\ref{tab: 单元格三种水平对齐方式}。


	若要调整具有多行内容的固定列宽单元格的垂直对齐方式，可以使用\shad{tabular}环境的\shad{p\{<宽度值>\}}、\shad{m\{<宽度值>\}}、\shad{b\{<宽度值>\}}参数，\textbf{它们分别表示使用对应单元格内容的上线、中线、底线作为垂直对齐时的基准线}，\textcolor{DarkRed}{而非使单元格内容上端对齐、垂直居中对齐、下端对齐}，效果见表\ref{tab: 单元格三种垂直对齐方式}。如果不好理解，不妨看看这个：\href{https://tex.stackexchange.com/questions/35293}{\ttfamily\color{DarkRed}\CJKunderline*[thickness=0.5bp, format=\color{DarkRed}]{p, m and b columns in tables}}。
	
	\begin{table}[!ht]
		\caption{\shad{单元格三种水平对齐方式}} \label{tab: 单元格三种水平对齐方式}
		\begin{tabular}{p{2cm} p{3cm}<{\centering} p{7cm}<{\raggedleft}}
			\toprule
			\textbf{姓名} & \textbf{所属势力} & \textbf{武功绝学} \\
			\midrule
			郭靖 & 重阳宫 & 降龙十八掌 \\
			黄蓉 & 丐帮 & 打狗棒法 \\
			洪七公 & 丐帮 & 降龙十八掌、打狗棒法 \\
			黄老邪 & 桃花岛 & 弹指神通、落英神剑掌、玉箫剑法 \\
			老顽童 & 重阳宫 & 左右互博术 \\
			一灯 & 云南大理 & 一阳指、千里传音 \\
			\bottomrule
		\end{tabular}
	\end{table}


	\begin{table}[!ht]
		\caption{\shad{单元格三种垂直对齐方式}} \label{tab: 单元格三种垂直对齐方式}
		\begin{tabular}{p{2cm} b{2cm} m{2cm}}
			\toprule
			\textbf{姓名} & \textbf{所属势力} & \textbf{武功绝学} \\
			\midrule
			洪七公 & 丐帮 & 降龙十八掌、打狗棒法 \\
			黄老邪黄老邪黄老邪 & 桃花岛桃花岛桃花岛 & 弹指神通、落英神剑掌、玉箫剑法 \\
			\bottomrule
		\end{tabular}
	\end{table}
	
	\newpage
	\subsubsection{带附注表格}
	
	更需要说明的是生成带附注的表格。本模板采用\shad{threeparttable}宏包实现将表格中的附注内容顶格排版在表格底部：
	\begin{enumerate}
		\item 使用\shadcmd{tnote\{<label>\}}在表格中插入上标编号；
		\setcounter{enumi}{98}% 更改后续列表编号的起始值
		\item 使用\shad{tablenotes}环境在表格底部排版附注。该环境提供选项\shad{online}用于将附注文本前的标号从默认的上标样式（见表\ref{tab: 江湖势力背调}）更改为非上标样式（见表\ref{tab: 已习得武功}）。
	\end{enumerate}
	
	\begin{table}[!ht]
		\caption{江湖势力背调} \label{tab: 江湖势力背调}
		\begin{threeparttable}
			\begin{tabular}{p{2cm} p{3cm}<{\centering} p{7cm}<{\raggedleft}}
				\toprule
				\textbf{姓名} & \textbf{所属势力} & \textbf{武功绝学} \\
				\midrule
				郭靖 & 重阳宫 & 降龙十八掌 \\
				黄蓉 & 丐帮 & 打狗棒法 \\
				洪七公 & 丐帮 & 降龙十八掌、打狗棒法 \\
				黄老邪 & 桃花岛 & 弹指神通、落英神剑掌、玉箫剑法 \\
				老顽童 & 重阳宫 & 左右互博术\tnote{1} \\
				一灯 & 云南大理 & 一阳指\tnote{2}、千里传音 \\
				\bottomrule
			\end{tabular}
			\begin{tablenotes}
				\item[1] 左右互搏术是金庸小说《射雕英雄传》中「老顽童」周伯通在桃花岛的地洞中创出的武功，本质是一心二用，能够两手同时做不同的事情，在金庸武侠体系中是一门非常精妙的武学，其对于人物本身的战斗力加成堪称台阶性。
				\item[2] 云南大理段氏嫡传的武功，在点穴功夫中位居天下第一，运功后以右手食指点穴，出指可缓可快，缓时潇洒飘逸，快则疾如闪电，但着指之处，分毫不差。当与敌挣搏凶险之际，用此指法既可贴近径点敌人穴道，也可从远处欺近身去，一中即离，一攻而退，实为克敌保身的无上妙术。
			\end{tablenotes}
		\end{threeparttable}
	\end{table}
	
	\begin{table}[!ht]
		\caption{已习得武功} \label{tab: 已习得武功}
		\begin{threeparttable}
			\begin{tabular}{p{3cm} p{3cm} p{5cm}}
				\toprule
				\textbf{武功绝学} & \textbf{传授者} & \textbf{传授地点} \\
				\midrule
				蛤蟆功 & 欧阳锋 & 重阳山脉 \\
				九阴真经 & 小龙女 & 活死人墓 \\
				打狗棒法 & 洪七公、黄蓉 & 华山之巅、英雄大会 \\
				玉箫剑法 & 黄老邪 & 深山老林 \\
				黯然销魂掌\tnote{1} & 自创 & 海边 \\
				\bottomrule
			\end{tabular}
			\begin{tablenotes}[online]
				\item[1] 黯然销魂掌，是在杨过与小龙女离别后，认为今生再也见不到小龙女，悲从中来，由此创作了黯然销魂掌。黯然销魂掌和心情有关，此后杨过与小龙女重逢后，其心理愉悦，故使不出黯然销魂掌。
			\end{tablenotes}
		\end{threeparttable}
	\end{table}

	上述方式排版的带附注表格无法通过点击表格中的编号跳转到对应附注。为此，本模板提供命令\shadcmd{puttablenotelabel\{<标签>\}}和\shadcmd{tablenoteref\{<标签>\}}来实现该操作：

	\begin{itemize}
		\item 首先将\shad{tablenotes}环境中手动设置的编号替换为\shadcmd{puttablenotelabel\{<标签>\}}；
		\item 然后在表格内容的对应位置使用\shadcmd{tablenoteref\{<标签>\}}即可。编号将自动生成，并按照使用\shadcmd{puttablenotelabel}的顺序递加。
	\end{itemize}

	这种方式由于需要建立交叉引用，通常需要用户编译两次。切记，\textbf{标签必须全文唯一}。表\ref{tab: 江湖势力背调（基于puttablenotelabel和tablenoteref）}提供了使用示例。

	\begin{table}[!ht]
		\caption{江湖势力背调（基于\shadcmd{puttablenotelabel}和\shadcmd{tablenoteref}）} \label{tab: 江湖势力背调（基于puttablenotelabel和tablenoteref）}
		\begin{threeparttable}
			\begin{tabular}{p{2cm} p{3cm} p{7cm}}
				\toprule
				\textbf{姓名} & \textbf{所属势力} & \textbf{武功绝学} \\
				\midrule
				郭靖 & 重阳宫 & 降龙十八掌 \\
				黄蓉 & 丐帮 & 打狗棒法 \\
				洪七公 & 丐帮 & 降龙十八掌、打狗棒法 \\
				黄老邪 & 桃花岛 & 弹指神通、落英神剑掌、玉箫剑法 \\
				老顽童 & 重阳宫 & 左右互博术\tablenoteref{tn: 左右互搏术} \\
				一灯 & 云南大理 & 一阳指\tablenoteref{tn: 一阳指}、千里传音 \\
				\bottomrule
			\end{tabular}
			\begin{tablenotes}
				\item[\puttablenotelabel{tn: 左右互搏术}] 左右互搏术是金庸小说《射雕英雄传》中「老顽童」周伯通在桃花岛的地洞中创出的武功，本质是一心二用，能够两手同时做不同的事情，在金庸武侠体系中是一门非常精妙的武学，其对于人物本身的战斗力加成堪称台阶性。
				\item[\puttablenotelabel{tn: 一阳指}] 云南大理段氏嫡传的武功，在点穴功夫中位居天下第一，运功后以右手食指点穴，出指可缓可快，缓时潇洒飘逸，快则疾如闪电，但着指之处，分毫不差。当与敌挣搏凶险之际，用此指法既可贴近径点敌人穴道，也可从远处欺近身去，一中即离，一攻而退，实为克敌保身的无上妙术。
			\end{tablenotes}
		\end{threeparttable}
	\end{table}
	
	
	\clearpage
	\subsubsection{跨页表格}
	
	原则上，长度不足一页的表格不应跨页。而对于本身超过一页的表格，本模板使用\href{https://mirrors.tuna.tsinghua.edu.cn/CTAN/macros/latex/required/tools/longtable.pdf}{\ttfamily\color{DarkRed}\CJKunderline*[thickness=0.5bp, format=\color{DarkRed}]{longtable}}宏包提供的\shad{longtable}环境实现。用户需要了解\shad{longtable}环境的基本使用方法，它与\shad{tabular}环境的最大区别在于需要用户自行定义分页后的表题、表头以及表尾。
	
	本模板提供了命令\shadcmd{CPcaption\{<当前表格总列数>\}\{<跨页表题>\}}来排版跨页之后的表题。\textbf{务必使用此命令}，否则，跨页表题的段前端后间距将不协调，且在其长度超过表宽时，无法获得预期的排版结果。
	
	此外，\textbf{\textcolor{DarkRed}{不应将\shad{longtable}环境嵌套在\shad{table}等浮动环境中}}，否则长表格将无法正常跨页。具体细节参见本小节示例表\ref{tab: 中国计算机学会部分推荐期刊及会议}和表\ref{tab: 中国计算机学会部分推荐期刊及会议简表}。
	
	% \newpage

	\begin{longtable}{p{2em} p{4.5em} p{11em} p{6em} p{11em}}
		\caption{中国计算机学会部分推荐期刊及会议} \label{tab: 中国计算机学会部分推荐期刊及会议} \\
		
		\toprule
		\textbf{序号} & \textbf{刊物简称} & \textbf{刊物全称} & \textbf{出版社} & \textbf{网址} \\
		\midrule
		\endfirsthead
		
		% 在这里设计首页以外的表题和表头
		\CPcaption{5}{中国计算机学会部分推荐期刊及会议}\\
		\toprule
		\textbf{序号} & \textbf{刊物简称} & \textbf{刊物全称} & \textbf{出版社} & \textbf{网址} \\
		\midrule
		\endhead
		
		% 在这里设计首页以外的表尾
		\bottomrule
		\multicolumn{5}{l}{续下页} \\  % 如不希望跨页表尾显示任何内容则注释掉即可
		\endfoot
		
		\bottomrule
		\endlastfoot
		
		1 & JSAC & IEEE Journal on Selected Areas in Communications & IEEE & http://dblp.uni-trier.de/db/journals/jsac/ \\
		2 & TMC & IEEE Transactions on Mobile Computing & IEEE & http://dblp.uni-trier.de/db/journals/tmc/ \\
		3 & TON & IEEE/ACM Transactions on Networking & IEEE/ACM & http://dblp.uni-trier.de/db/journals/ton/ \\
		1 & TOIT & ACM Transactions on Internet Technology & ACM & http://dblp.uni-trier.de/db/journals/toit/ \\
		2 & TOMM & ACM Transactions on Multimedia Computing, Communications and Applications & ACM & http://dblp.uni-trier.de/db/journals/tomccap/ \\
		3 & TOSN & ACM Transactions on Sensor Networks & ACM & http://dblp.uni-trier.de/db/journals/tosn/ \\
		4 & CN & Computer Networks & Elsevier & http://dblp.uni-trier.de/db/journals/cn/ \\
		5 & TCOM & IEEE Transactions on Communications & IEEE & http://dblp.uni-trier.de/db/journals/tcom/ \\
		6 & TWC & IEEE Transactions on Wireless Communications & IEEE & http://dblp.uni-trier.de/db/journals/twc/ \\
		1 & & Ad Hoc Networks & Elsevier & http://dblp.uni-trier.de/db/journals/adhoc/ \\
		2 & CC & Computer Communications & Elsevier & http://dblp.uni-trier.de/db/journals/comcom/ \\
		3 & TNSM & IEEE Transactions on Network and Service Management & IEEE & http://dblp.uni-trier.de/db/journals/tnsm/ \\
		4 & & IET Communications & IET & http://dblp.uni-trier.de/db/journals/iet-com/ \\
		5 & JNCA & Journal of Network and Computer Applications & Elsevier & http://dblp.uni-trier.de/db/journals/jnca/ \\
		6 & MONET & Mobile Networks and Applications & Springer & http://dblp.uni-trier.de/db/journals/monet/ \\
		7 & & Networks & Wiley & http://dblp.uni-trier.de/db/journals/networks/ \\
		8 & PPNA & Peer-to-Peer Networking and Applications & Springer & http://dblp.uni-trier.de/db/journals/ppna/ \\
		9 & WCMC & Wireless Communications and Mobile Computing & Wiley & http://dblp.uni-trier.de/db/journals/wicomm/ \\
		10 & & Wireless Networks & Springer & http://dblp.uni-trier.de/db/journals/winet/ \\
		11 & IOT & IEEE Internet of Things Journal & IEEE & https://dblp.org/db/journals/ iotj/index.html \\
		1 & SIGCOMM & ACM International Conference on Applications, Technologies, Architectures, and Protocols for Computer Communication & ACM & http://dblp.uni-trier.de/db/conf/sigcomm/ index.html \\
		2 & MobiCom & ACM International Conference on Mobile Computing and Networking & ACM & http://dblp.uni-trier.de/db/conf/mobicom/ \\
		3 & INFOCOM & IEEE International Conference on Computer Communications & IEEE & http://dblp.uni-trier.de/db/conf/infocom/ \\
		4 & NSDI & Symposium on Network System Design and Implementation & USENIX & http://dblp.uni-trier.de/db/conf/nsdi/ \\
		1 & SenSys & ACM Conference on Embedded Networked Sensor Systems & ACM & http://dblp.uni-trier.de/db/conf/sensys/ \\
		2 & CoNEXT & ACM International Conference on Emerging Networking Experiments and Technologies & ACM & http://dblp.uni-trier.de/db/conf/conext/ \\
		3 & SECON & IEEE International Conference on Sensing, Communication, and Networking & IEEE & http://dblp.uni-trier.de/db/conf/secon/ \\
		4 & IPSN & International Conference on Information Processing in Sensor Networks & IEEE/ACM & http://dblp.uni-trier.de/db/conf/ipsn/ \\
		5 & MobiSys & ACM International Conference on Mobile Systems, Applications, and Services & ACM & http://dblp.uni-trier.de/db/conf/mobisys/ \\
		6 & ICNP & IEEE International Conference on Network Protocols & IEEE & http://dblp.uni-trier.de/db/conf/icnp/ \\
		7 & MobiHoc & International Symposium on Theory, Algorithmic Foundations, and Protocol Design for Mobile Networks and Mobile Computing & ACM/IEEE & http://dblp.uni-trier.de/db/conf/mobihoc/ \\
		8 & NOSSDAV & International Workshop on Network and Operating System Support for Digital Audio and Video & ACM & http://dblp.uni-trier.de/db/conf/nossdav/ \\
		9 & IWQoS & IEEE/ACM International Workshop on Quality of Service & IEEE & http://dblp.uni-trier.de/db/conf/iwqos/ \\
		10 & IMC & ACM Internet Measurement Conference & ACM/USENIX & http://dblp.uni-trier.de/db/conf/imc/ \\
		
		
	\end{longtable}

	\newpage

	\begin{longtable}{p{2em} p{4.5em}}
		\caption{中国计算机学会部分推荐期刊及会议简表（用于测试跨页表宽低于表题长度的情况）} \label{tab: 中国计算机学会部分推荐期刊及会议简表} \\
		
		\toprule
		\textbf{序号} & \textbf{刊物简称} \\
		\midrule
		\endfirsthead
		
		% 在这里设计首页以外的表题和表头
		\CPcaption{2}{中国计算机学会部分推荐期刊及会议简表（用于测试跨页表宽低于表题长度的情况）}\\
		\toprule
		\textbf{序号} & \textbf{刊物简称} \\
		\midrule
		\endhead
		
		% 在这里设计首页以外的表尾
		\bottomrule
		\multicolumn{2}{r}{续下页} \\  % 如不希望跨页表尾显示任何内容则注释掉即可
		\endfoot
		
		\bottomrule
		\endlastfoot
		
		1 & JSAC \\
		2 & TMC \\
		3 & TON \\
		1 & TOIT \\
		2 & TOMM \\
		3 & TOSN \\
		4 & CN \\
		5 & TCOM \\
		6 & TWC \\
		2 & CC \\
		3 & TNSM \\
		5 & JNCA \\
		6 & MONET \\
		8 & PPNA \\
		9 & WCMC \\
		11 & IOT \\
		1 & SIGCOMM \\
		2 & MobiCom \\
		3 & INFOCOM \\
		4 & NSDI \\
		1 & SenSys \\
		2 & CoNEXT \\
		3 & SECON \\
		4 & IPSN \\
		5 & MobiSys \\
		6 & ICNP \\
		7 & MobiHoc \\
		8 & NOSSDAV \\
		9 & IWQoS \\
		10 & IMC \\
		1 & JSAC \\
		2 & TMC \\
		3 & TON \\
		1 & TOIT \\
		2 & TOMM \\
		3 & TOSN \\
		4 & CN \\
		5 & TCOM \\
		6 & TWC \\
		2 & CC \\
		3 & TNSM \\
		5 & JNCA \\
		6 & MONET \\
		8 & PPNA \\
		9 & WCMC \\
		11 & IOT \\
		1 & SIGCOMM \\
		2 & MobiCom \\
		3 & INFOCOM \\
		4 & NSDI \\
		1 & SenSys \\
		2 & CoNEXT \\
		3 & SECON \\
		4 & IPSN \\
		5 & MobiSys \\
		6 & ICNP \\
		7 & MobiHoc \\
		8 & NOSSDAV \\
		9 & IWQoS \\
		10 & IMC \\
	\end{longtable}

	\clearpage
	\subsubsection{跨页带附注表格}

	很遗憾，\shad{threeparttable}无法做到跨页，而\shad{longtable}又无法像前者那样稳定完美地排版附注。无奈之下，本模板提供了一种繁琐但可行的做法。

	思路是利用\shad{longtable}提供的表尾自定义功能来插入附注，即用户需要修改\shad{longtable}在\shadcmd{endlastfoot}前对表格尾部的设置。为此，模板提供命令：
	
	\shadcmd{tablenotetext(online)(<附注编号悬挂距离>)[<附注上方垂直间隔>]\{<附注总宽度>\}\{<附注标签>\}\{<附注内容>\}}。

	\begin{itemize}
		\item 此命令的第一项可选参数以\shad{()}标识，它仅接受参数\shad{online}，作用是将附注编号从默认的上标形式更改为行内形式，即模仿\shad{tablenotes}环境的选项。此间差异可参考表\ref{tab: 跨页带附注表格示例}的附注；
		\item 第二项可选参数仅在设置了第一项可选参数为\shad{online}时有效，用于调整附注编号的悬挂缩进距离，默认是\shad{1em}。当表格的附注数量过多时，可通过手动调整该可选参数来避免过长的编号与后方文字重叠，见表\ref{tab: 跨页带附注表格恰巧在附注开始处换页示例}；
		\item 第三项可选参数以\shad{[]}标识，在生成的附注与上方内容间的垂直距离不合适时，可通过此可选参数手动对其进行调整，默认为\shad{0bp}；
		\item 第四项强制参数对应附注整体的宽度，其取值需要根据表格排版后的宽度反复调整，直至附注与表格尾线等宽为止。这的确很繁琐，但我实在能力有限，想不出更好的办法，自动确定表格的实际宽度真的很难；
		\item 第五项强制参数负责设置附注编号对应的标签，以便后续在表格中使用\shadcmd{tablenoteref\{\}}进行引用，标签同样需要全文唯一；
		\item 第六项强制参数对应附注的实际内容。
	\end{itemize}

	\textcolor{red}{\textbf{注意：}}使用\shadcmd{tablenotetext}命令的前提是对表格首列进行特定设置，用户\textbf{必须}通过\shad{p\{<长度>\}}或\shad{m\{<长度>\}}来人为指定\shad{longtable}的首列宽度，\textcolor{red}{\textbf{切不可}}将之设置为\shad{c}、\shad{r}或\shad{l}。另外，表格中最后一条\shadcmd{tablenotetext}之后不需要\shadcmd{\shadcmd{}}，否则表尾与下方文本的间隔将偏大。

	可能出现一种极端情况：跨页表格恰好在附注开始处分页，此时也会产生另类的排版结果。要解决该问题，需手动在表格内容的适当位置使用\href{https://mirrors.tuna.tsinghua.edu.cn/CTAN/macros/latex/required/tools/longtable.pdf}{\ttfamily\color{DarkRed}\CJKunderline*[thickness=0.5bp, format=\color{DarkRed}]{longtable}}宏包提供的\shadcmd{pagebreak}命令提前断页，代价是前一页底部可能有更多空白，参考表\ref{tab: 跨页带附注表格恰巧在附注开始处换页示例}中的做法。

	\shadcmd{tablenotetext}命令不算稳定。其参数在不同的表格中需要特调，甚至在同一表格的不同排版设置下都是如此，一表一参。若不得不用，还望仔细。


	\newpage

	\begin{longtable}{p{2em} p{4.5em} p{20em} p{6em}}
		\caption{跨页带附注表格示例} \label{tab: 跨页带附注表格示例} \\
		
		\toprule
		\textbf{序号} & \textbf{刊物简称} & \textbf{刊物全称} & \textbf{出版社} \\
		\midrule
		\endfirsthead
		
		% 在这里设计首页以外的表题和表头
		\CPcaption{4}{跨页带附注表格示例}\\
		\toprule
		\textbf{序号} & \textbf{刊物简称} & \textbf{刊物全称} & \textbf{出版社} \\
		\midrule
		\endhead
		
		% 在这里设计首页以外的表尾
		\bottomrule
		\multicolumn{4}{l}{续下页} \\  % 如不希望跨页表尾显示任何内容则注释掉即可
		\endfoot
		
		\bottomrule
		\tablenotetext[-7bp]{37.2em}{tn: 手动跨页表格附注标签IEEE}{IEEE是指电气和电子工程师学会，是一个国际性的专业学会，以促进电气工程、电子工程、计算机科学和相关领域的科学和技术发展为宗旨。成立于1884年，总部位于美国纽约。IEEE 的会员包括来自世界各地的专业人士、工程师、学者和学生，是全球最大的技术专业组织之一。} \\
		\tablenotetext(online)[6bp]{37.2em}{tn: 手动跨页表格附注标签ACM}{ACM是指国际计算机学会，成立于1947年，是一个国际性的科技教育组织，是世界上第一个科学性及教育性计算机学会，总部设在美国纽约。国际计算机学会是世界上最大的计算机领域专业性学术组织，汇集了国际计算机领域教育家，研究人员，工业界人士及学生。ACM致力于提高在中国的活动的规格与影响力。在此基础上，学会成立了ACM中国理事会，为在中国的学会会员与学会活动提供支持。}% 注意这里不需要\\
		\endlastfoot
		
		1 & JSAC & IEEE Journal on Selected Areas in Communications & IEEE \\
		2 & TMC & IEEE Transactions on Mobile Computing & IEEE \\
		3 & TON & IEEE/ACM Transactions on Networking & IEEE/ACM \\
		1 & TOIT & ACM Transactions on Internet Technology & ACM \\
		2 & TOMM & ACM Transactions on Multimedia Computing, Communications and Applications & ACM \\
		3 & TOSN & ACM Transactions on Sensor Networks & ACM \\
		4 & CN & Computer Networks & Elsevier \\
		5 & TCOM & IEEE Transactions on Communications & IEEE \\
		6 & TWC & IEEE Transactions on Wireless Communications & IEEE \\
		1 & & Ad Hoc Networks & Elsevier \\
		2 & CC & Computer Communications & Elsevier \\
		3 & TNSM & IEEE Transactions on Network and Service Management & IEEE \\
		4 & & IET Communications & IET \\
		5 & JNCA & Journal of Network and Computer Applications & Elsevier \\
		6 & MONET & Mobile Networks and Applications & Springer \\
		7 & & Networks & Wiley \\
		8 & PPNA & Peer-to-Peer Networking and Applications & Springer \\
		9 & WCMC & Wireless Communications and Mobile Computing & Wiley \\
		10 & & Wireless Networks & Springer \\
		11 & IOT & IEEE Internet of Things Journal & IEEE \\
		1 & SIGCOMM & ACM International Conference on Applications, Technologies, Architectures, and Protocols for Computer Communication & ACM \\
		2 & MobiCom & ACM International Conference on Mobile Computing and Networking & ACM \\
		3 & INFOCOM & IEEE International Conference on Computer Communications & IEEE \\
		4 & NSDI & Symposium on Network System Design and Implementation & USENIX \\
		1 & SenSys & ACM Conference on Embedded Networked Sensor Systems & ACM \\
		2 & CoNEXT & ACM International Conference on Emerging Networking Experiments and Technologies & ACM \\
		3 & SECON & IEEE International Conference on Sensing, Communication, and Networking & IEEE \\
		4 & IPSN & International Conference on Information Processing in Sensor Networks & IEEE/ACM \\
		5 & MobiSys & ACM International Conference on Mobile Systems, Applications, and Services & ACM \\
		6 & ICNP & IEEE International Conference on Network Protocols & IEEE\tablenoteref{tn: 手动跨页表格附注标签IEEE} \\
		7 & MobiHoc & International Symposium on Theory, Algorithmic Foundations, and Protocol Design for Mobile Networks and Mobile Computing & ACM/IEEE \\
		8 & NOSSDAV & International Workshop on Network and Operating System Support for Digital Audio and Video & ACM\tablenoteref{tn: 手动跨页表格附注标签ACM} \\
		9 & IWQoS & IEEE/ACM International Workshop on Quality of Service & IEEE \\
		10 & IMC & ACM Internet Measurement Conference & ACM/USENIX \\
		
		
	\end{longtable}

	表格后参照文本表格后参照文本表格后参照文本表格后参照文本表格后参照文本表格后参照文本表格后参照文本表格后参照文本表格后参照文本表格后参照文本表格后参照文本表格后参照文本表格后参照文本表格后参照文本表格后参照文本表格后参照文本表格后参照文本表格后参照文本表格后参照文本表格后参照文本表格后参照文本

	\newpage

	\begin{longtable}{m{2em}<{\centering} p{4.5em} p{15em} p{6em}}
		\caption{跨页带附注表格恰巧在附注开始处换页示例} \label{tab: 跨页带附注表格恰巧在附注开始处换页示例} \\
		
		\toprule
		\textbf{序号} & \textbf{刊物简称} & \textbf{刊物全称} & \textbf{出版社} \\
		\midrule
		\endfirsthead
		
		% 在这里设计首页以外的表题和表头
		\CPcaption{4}{跨页带附注表格恰巧在附注中需要换页示例}\\
		\toprule
		\textbf{序号} & \textbf{刊物简称} & \textbf{刊物全称} & \textbf{出版社} \\
		\midrule
		\endhead
		
		% 在这里设计首页以外的表尾
		\bottomrule
		\multicolumn{4}{l}{续下页} \\  % 如不希望跨页表尾显示任何内容则注释掉即可
		\endfoot
		
		\bottomrule
		\tablenotetext[5bp]{32.1em}{tn: 手动跨页表格附注标签Elsevier}{爱思唯尔，创办于1880年，属于RELX集团旗下，总部位于阿姆斯特丹。爱思唯尔是一家荷兰的国际化多媒体出版集团，主要为科学家、研究人员、学生、医学以及信息处理的专业人士提供信息产品和革新性工具。爱思唯尔是全球领先的科学与医学信息服务机构，旗下出版《柳叶刀》《细胞》等2800多种学术期刊。} \setcounter{tablenote}{99}\\
		\tablenotetext(online)(2em)[5bp]{32.1em}{tn: 手动跨页表格附注标签USENIX}{USENIX成立于1975年，当时的名字叫做Unix用户群。它的主要目的是学习及开发Unix以及类似系统。1977年六月，美国电话电报公司的律师告诉用户群他们不能继续使用UNIX这个名字，因为UNIX是美国电话电报公司所拥有的一个商标。所以这个用户群更名成USENIX.从那以后，USENIX逐渐发展成一个倍受尊敬的由计算机操作系统用户，开发者和研究者所组成的机构。}% 注意这里不需要\\
		\endlastfoot
		
		1 & JSAC & IEEE Journal on Selected Areas in Communications & IEEE \\
		2 & TMC & IEEE Transactions on Mobile Computing & IEEE \\
		3 & TON & IEEE/ACM Transactions on Networking & IEEE/ACM \\
		1 & TOIT & ACM Transactions on Internet Technology & ACM \\
		2 & TOMM & ACM Transactions on Multimedia Computing, Communications and Applications & ACM \\
		3 & TOSN & ACM Transactions on Sensor Networks & ACM \\
		4 & CN & Computer Networks & Elsevier \\
		5 & TCOM & IEEE Transactions on Communications & IEEE \\
		6 & TWC & IEEE Transactions on Wireless Communications & IEEE \\
		1 & & Ad Hoc Networks & Elsevier \\
		2 & CC & Computer Communications & Elsevier\tablenoteref{tn: 手动跨页表格附注标签Elsevier} \\
		3 & TNSM & IEEE Transactions on Network and Service Management & IEEE \\
		% 4 & & IET Communications & IET \\
		% 5 & JNCA & Journal of Network and Computer Applications & Elsevier \\
		% 6 & MONET & Mobile Networks and Applications & Springer \\
		% 7 & & Networks & Wiley \\
		% 8 & PPNA & Peer-to-Peer Networking and Applications & Springer \\
		% 9 & WCMC & Wireless Communications and Mobile Computing & Wiley \\
		% 10 & & Wireless Networks & Springer \\
		% 11 & IOT & IEEE Internet of Things Journal & IEEE \\
		1 & SIGCOMM & ACM International Conference on Applications, Technologies, Architectures, and Protocols for Computer Communication & ACM \\
		2 & MobiCom & ACM International Conference on Mobile Computing and Networking & ACM \\
		3 & INFOCOM & IEEE International Conference on Computer Communications & IEEE \\
		4 & NSDI & Symposium on Network System Design and Implementation & USENIX\tablenoteref{tn: 手动跨页表格附注标签USENIX} \\
		1 & SenSys & ACM Conference on Embedded Networked Sensor Systems & ACM \\
		\pagebreak  % 用户可以尝试注释掉这条命令看看现象
		2 & CoNEXT & ACM International Conference on Emerging Networking Experiments and Technologies & ACM \\
		% 3 & SECON & IEEE International Conference on Sensing, Communication, and Networking & IEEE \\
		% 4 & IPSN & International Conference on Information Processing in Sensor Networks & IEEE/ACM \\
		% 5 & MobiSys & ACM International Conference on Mobile Systems, Applications, and Services & ACM \\
		% 6 & ICNP & IEEE International Conference on Network Protocols & IEEE \\
		% 7 & MobiHoc & International Symposium on Theory, Algorithmic Foundations, and Protocol Design for Mobile Networks and Mobile Computing & ACM/IEEE \\
		% 8 & NOSSDAV & International Workshop on Network and Operating System Support for Digital Audio and Video & ACM \\
		% 9 & IWQoS & IEEE/ACM International Workshop on Quality of Service & IEEE \\
		% 10 & IMC & ACM Internet Measurement Conference & ACM/USENIX \\
		
		
		
	\end{longtable}

	表格后参照文本表格后参照文本表格后参照文本表格后参照文本表格后参照文本表格后参照文本表格后参照文本表格后参照文本表格后参照文本表格后参照文本表格后参照文本表格后参照文本表格后参照文本表格后参照文本表格后参照文本表格后参照文本表格后参照文本表格后参照文本表格后参照文本表格后参照文本表格后参照文本
	
	\clearpage
	\subsection{伪代码}
	
	模板基于\href{https://mirrors.sustech.edu.cn/CTAN/macros/latex/contrib/algorithm2e/doc/algorithm2e.pdf}{\ttfamily\color{DarkRed}\CJKunderline*[thickness=0.5bp, format=\color{DarkRed}]{algorithm2e}}宏包提供的\shad{algorithm}环境排版伪代码，默认不添加左右侧框线，且顶部框线和底部框线仿照表格进行了加粗，字体大小也调整到了\textbf{五号字}，与表格保持一致。若要恢复伪代码的左右侧框线，可以在载入文档类时使用\shad{boxruled}选项。该环境生成的伪码与正文文本保持相同宽度。
	
	除了\href{https://mirrors.sustech.edu.cn/CTAN/macros/latex/contrib/algorithm2e/doc/algorithm2e.pdf}{\ttfamily\color{DarkRed}\CJKunderline*[thickness=0.5bp, format=\color{DarkRed}]{algorithm2e}}宏包本身提供的各种条件、循环语句，本模板基于宏包提供的接口，追加了\shad{Do While}和\shad{Loop}循环语句：
	\begin{itemize}
		\item \shadcmd{DoWhile(<紧跟关键字do的文本，可用于添加注释>)\{<循环条件>\}\{<循环体>\}}
		\item \shadcmd{Loop(<紧跟关键字loop的文本，可用于添加注释>)\{<循环体>\}}
	\end{itemize}
	
	
	此外，基于调整后的\shad{algorithm2e}环境，本模板进一步封装了\shad{algo}环境，它将生成比\shad{algorithm}环境更\textbf{\color{DarkRed}窄}的伪码浮动区域。除了接受浮动可选参数\shad{[htbp]}，\shad{algo}环境还支持另一可选参数\shad{(<伪码距正文文本边界的总距离>)}：该参数控制浮动体距正文文本边界的总距离，默认\shad{4em}，即单边缩进\shad{2em}，与其下首行文本对齐。两可选参数可单独或同时使用，\textbf{同时使用的顺序必须与下方示例保持一致：}
	
	\begin{verbatim}
		\begin{algo}[<浮动选项>](<伪码距正文文本边界的总距离>)
		    .....
		\end{algo}
	\end{verbatim}
	
	算法\ref{alg: algorithm环境伪码示例}和算法\ref{alg: algo环境伪码示例}分别展示了两种环境默认生成的伪码样式；过程\ref{alg: algorithm环境修改伪码标签示例}和过程\ref{alg: algo环境修改伪码标签并调整宽度示例}展示了如何修改伪码中的一些标签，以及如何调整\shad{algo}生成的伪码宽度。

	
	\begin{algorithm}[!h]
		\caption{algorithm环境伪码示例} \label{alg: algorithm环境伪码示例}
		\Input{1) 输入1;\newline 2) 输入2。}
		\Output{输出结果。}
		伪码行1。
		
		\For(\tcc*[f]{循环条件注释1}){循环条件1}{
			伪码行2。
			
			\tcp{注释2}
			伪码行3。
			
			\DoWhile(\tcc*[f]{循环条件注释3}){循环条件2}{
				伪码行4。
			}
			
			\tcc{loop循环}
			\Loop(\tcc*[f]{注释4}){
				循环体1。
			}
			
			\Repeat(\tcc*[f]{循环条件注释5}){循环条件3}{
				循环体2。
			}
			
			\tcp{if-elseif-else结构示例}
			\uIf(\tcc*[f]{条件注释6}){条件语句5}{
				条件语句5为真，伪码行5。
			}
			\uElseIf(\tcc*[f]{elseif条件语句}){条件语句6}{
				条件语句6为真，伪码行6。
			}
			\Else{
				条件5和6均为假，伪码行7。\tcp*[f]{else代码内容}
			}
			
			\If(\tcc*[f]{条件注释7}){条件语句7}{
				伪码行8。
			}
		}
		\textbf{return} 算法结果。
	\end{algorithm}
	
	\begin{algorithm}[!h]
		\renewcommand{\algorithmcfname}{过程}  % 修改伪码标签需要在\caption{}之前
		\caption{algorithm环境临时修改伪码标签示例} \label{alg: algorithm环境修改伪码标签示例}
		\SetKwInOut{Input}{In}
		\SetKwInOut{Output}{Out}
		\Input{1) 输入1; 2) 输入2。}
		\Output{输出结果。}
		伪码行1。
		
		\For(\tcc*[f]{循环条件注释1}){循环条件1}{
			伪码行2。
			
			\tcp{注释2}
			伪码行3。
			
			\DoWhile(\tcc*[f]{循环条件注释3}){循环条件2}{
				伪码行4。
			}
			
			\tcc{loop循环}
			\Loop(\tcc*[f]{注释4}){
				循环体1。
			}
			
			\Repeat(\tcc*[f]{循环条件注释5}){循环条件3}{
				循环体2。
			}
			\eIf(\tcc*[f]{条件注释6}){条件语句6}{
				为真，伪码行5。
			}{
				条件为假，伪码行6。\tcp*[f]{else代码内容}
			}
			
			\If(\tcc*[f]{条件注释7}){条件语句7}{
				伪码行7。
			}
		}
		\textbf{return} 算法结果。
	\end{algorithm}
	
	\begin{algo}[!h]
		\caption{algo环境伪码示例} \label{alg: algo环境伪码示例}
		\Input{1) 输入1;\newline 2) 输入2。}
		\Output{输出结果。}
		伪码行1。
		
		\For(\tcc*[f]{循环条件注释1}){循环条件1}{
			伪码行2。
			
			\tcp{注释2}
			伪码行3。
			
			\DoWhile(\tcc*[f]{循环条件注释3}){循环条件2}{
				伪码行4。
			}
			
			\tcc{loop循环}
			\Loop(\tcc*[f]{注释4}){
				循环体1。
			}
			
			\Repeat(\tcc*[f]{循环条件注释5}){循环条件3}{
				循环体2。
			}
			\eIf(\tcc*[f]{条件注释6}){条件语句6}{
				为真，伪码行5。
			}{
				条件为假，伪码行6。\tcp*[f]{else代码内容}
			}
			
			\If(\tcc*[f]{条件注释7}){条件语句7}{
				伪码行7。
			}
		}
		\textbf{return} 算法结果。
	\end{algo}
	
	\begin{algo}[!h](8em)
		\renewcommand{\algorithmcfname}{过程}  % 修改伪码标签需要在\caption{}之前
		\caption{algo环境临时修改伪码标签并调整宽度示例} \label{alg: algo环境修改伪码标签并调整宽度示例}
		\SetKwInOut{Input}{In}
		\SetKwInOut{Output}{Out}
		\Input{1) 输入1; 2) 输入2。}
		\Output{输出结果。}
		伪码行1。
		
		\For(\tcc*[f]{循环条件注释1}){循环条件1}{
			伪码行2。
			
			\tcp{注释2}
			伪码行3。
			
			\DoWhile(\tcc*[f]{循环条件注释3}){循环条件2}{
				伪码行4。
			}
			
			\tcc{loop循环}
			\Loop(\tcc*[f]{注释4}){
				循环体1。
			}
			
			\Repeat(\tcc*[f]{循环条件注释5}){循环条件3}{
				循环体2。
			}
			\eIf(\tcc*[f]{条件注释6}){条件语句6}{
				为真，伪码行5。
			}{
				条件为假，伪码行6。\tcp*[f]{else代码内容}
			}
			
			\If(\tcc*[f]{条件注释7}){条件语句7}{
				伪码行7。
			}
		}
		\textbf{return} 算法结果。
	\end{algo}

	\clearpage
	\subsection{各种列表}

	本模板对\shad{itemize}、\shad{enumerate}和\shad{description}这三种基本列表环境进行了设置，用户可根据实际情况选用。

	\shad{itemize}示例：

	\begin{itemize}
		\item 外层列表条目1。多行填充多行填充多行填充多行填充多行填充多行填充多行填充多行填充多行填充多行填充
		\item 外层列表条目2
		\begin{itemize}
			\item 内层列表条目1。多行填充多行填充多行填充多行填充多行填充多行填充多行填充多行填充多行填充多行填充
			\item 内层列表条目2
		\end{itemize}
		\item 外层列表条目3
	\end{itemize}

	\null

	\shad{enumerate}示例：

	\begin{enumerate}
		\item 外层枚举条目1。多行填充多行填充多行填充多行填充多行填充多行填充多行填充多行填充多行填充多行填充
		\item 外层枚举条目2
		\begin{enumerate}
			\item 内层枚举条目1。多行填充多行填充多行填充多行填充多行填充多行填充多行填充多行填充多行填充多行填充
			\item 内层枚举条目2
		\end{enumerate}
		\item 外层枚举条目3
	\end{enumerate}

	\null

	\shad{description}示例:

	\begin{description}
		\item[描述1] 外层描述条目1。多行填充多行填充多行填充多行填充多行填充多行填充多行填充多行填充多行填充多行填充
		\item[描述2] 外层描述条目2
		\begin{description}
			\item[描述2.1] 内层描述条目1。多行填充多行填充多行填充多行填充多行填充多行填充多行填充多行填充多行填充多行填充
			\item[描述2.2] 内层描述条目2
		\end{description}
		\item[描述3] 外层描述条目3
	\end{description}
	
	\clearpage
	\subsection{定义、公理、定理、命题、推论、引理、示例、假设、注、证明}
	
	本模板分别定义了环境：\shad{definition}、\shad{axiom}、\shad{theorem}、\shad{proposition}、\shad{corollary}、\shad{lemma}、\shad{example}、\shad{assumption}、\shad{annotation}和\shad{proof}。示例如下：

	\begin{definition}[具体名称]
		云南大理段氏嫡传的武功，在点穴功夫中位居天下第一，运功后以右手食指点穴，出指可缓可快，缓时潇洒飘逸，快则疾如闪电。
	\end{definition}

	\begin{axiom}[具体名称]
		云南大理段氏嫡传的武功，在点穴功夫中位居天下第一，运功后以右手食指点穴，出指可缓可快，缓时潇洒飘逸，快则疾如闪电。
	\end{axiom}
	
	\begin{theorem}[具体名称]
		云南大理段氏嫡传的武功，在点穴功夫中位居天下第一，运功后以右手食指点穴，出指可缓可快，缓时潇洒飘逸，快则疾如闪电。

		\begin{enumerate}
			\item 当与敌挣搏凶险之际，用此指法既可贴近径点敌人穴道，也可从远处欺近身去，一中即离，一攻而退，实为克敌保身的无上妙术。
			\begin{enumerate}
				\item 当与敌挣搏凶险之际，用此指法既可贴近径点敌人穴道，也可从远处欺近身去，一中即离，一攻而退，实为克敌保身的无上妙术。
			\end{enumerate}
		\end{enumerate}
	\end{theorem}
	
	\begin{proposition}[具体名称]
		云南大理段氏嫡传的武功，在点穴功夫中位居天下第一，运功后以右手食指点穴，出指可缓可快，缓时潇洒飘逸，快则疾如闪电。
	\end{proposition}
	
	\begin{corollary}[具体名称]
		云南大理段氏嫡传的武功，在点穴功夫中位居天下第一，运功后以右手食指点穴，出指可缓可快，缓时潇洒飘逸，快则疾如闪电。
	\end{corollary}
	
	\begin{lemma}[具体名称]
		云南大理段氏嫡传的武功，在点穴功夫中位居天下第一，运功后以右手食指点穴，出指可缓可快，缓时潇洒飘逸，快则疾如闪电。
	\end{lemma}

	\begin{example}[具体名称]
		云南大理段氏嫡传的武功，在点穴功夫中位居天下第一，运功后以右手食指点穴，出指可缓可快，缓时潇洒飘逸，快则疾如闪电。
	\end{example}

	\begin{assumption}[具体名称]
		云南大理段氏嫡传的武功，在点穴功夫中位居天下第一，运功后以右手食指点穴，出指可缓可快，缓时潇洒飘逸，快则疾如闪电。
	\end{assumption}

	\begin{annotation}[具体名称]
		云南大理段氏嫡传的武功，在点穴功夫中位居天下第一，运功后以右手食指点穴，出指可缓可快，缓时潇洒飘逸，快则疾如闪电。
	\end{annotation}
	
	\begin{proof}
		云南大理段氏嫡传的武功，在点穴功夫中位居天下第一，运功后以右手食指点穴，出指可缓可快，缓时潇洒飘逸，快则疾如闪电。
		\begin{itemize}
			\item 当与敌挣搏凶险之际，用此指法既可贴近径点敌人穴道，也可从远处欺近身去，一中即离，一攻而退，实为克敌保身的无上妙术。
			\begin{itemize}
				\item 当与敌挣搏凶险之际，用此指法既可贴近径点敌人穴道，也可从远处欺近身去，一中即离，一攻而退，实为克敌保身的无上妙术。
			\end{itemize}
		\end{itemize}
	\end{proof}

	\subsection{脚注}
	
	本模板使用包含了带圈数字的字体来替换LaTeX绘制的带圈数字，提供了充足的带圈编号数量，同时保证了带圈脚注编号足够优雅。
	
	在正文中加入脚注直接在需要放置脚注标签的位置使用\shadcmd{footnote\{<脚注内容>\}}即可。
	
	在其他环境中，如表格，则需要需要使用\shadcmd{footnotemark}配合\shadcmd{footnotetext\{<脚注文本>\}}。在需要放置脚注标签的位置使用\shadcmd{footnotemark}，然后在环境外使用\shadcmd{footnotetext\{<脚注文本>\}}指明脚注内容\footnote{更详细的使用方法参考\href{https://blog.csdn.net/xovee/article/details/127563209}{\ttfamily\color{DarkRed}\CJKunderline*[thickness=0.5bp, format=\color{DarkRed}]{LaTeX脚注}}。冗余文本用于展示脚注内容发生换行后的情况；冗余文本用于展示脚注内容发生换行后的情况；冗余文本用于展示脚注内容发生换行后的情况。}。
	
	
	\subsection{模板中的各种编号}
	
	“标题”、“图片”、“表格”、“伪码”、“公式”、“定义”、“定理”、“命题”、“推论”、“引理”、“证明”、“脚注”这些文档元素均可自行计算并生成编号。\textbf{但形如\textcolor{DodgerBlue}{(1-1a)}的子公式编号不能完全自动生成}，为此，模板提供了\shadcmd{subeqtag[<子公式编号标签>]}命令。

	在为数学模型的约束创建编号时，常见的方式可能是使用\shadcmd{tag\{\}}命令直接指定编号内容。但是，该方式操作繁琐，且在后续需要调换或增删约束时很容易漏改某些tag，导致子公式编号混乱，而又不容易察觉。

	本模板提供的\shadcmd{subeqtag[<子公式编号标签>]}命令彻底避免了上述问题。使用者只需要在对应的约束后插入\shadcmd{subeqtag}，即可赋予该约束与当前主公式编号保持一致的次级编号。而且，对连续多个约束使用该命令会\textbf{自动生成}递增的子公式编号，交换约束顺序编号也会自行更新，断不会出错。
	
	如果需要在正文中引用某个子公式编号，那么可以像往常一样在\shadcmd{subeqtag}之后使用\shadcmd{label\{<编号标签>\}}，或者直接指定\shadcmd{subeqtag[<子公式编号标签>]}的可选参数，非常人性化。
	
	下面的源码将产生式\eqref{eq: obj 1}\textasciitilde \eqref{eq: constriant gamma}对应的例子，其中式\eqref{eq: constraint x}和\eqref{eq: constriant gamma}使用了\shadcmd{subeqtag[<子公式编号标签>]}的可选参数。
	
	
	\begin{verbatim}
		\begin{align}
			\max \log \left(x^2 + y^2 + z^2 + v^2 + g^2 + m^2 + k^2\right)
			\label{eq: obj 1} \\
			\text{s.t.} \quad x \leq 1, \subeqtag[eq: constraint x] \\
			y \leq 2, \subeqtag \\
			z \leq 4, \subeqtag
		\end{align}
		
		\begin{align}
			\min \left(\boldsymbol{\alpha} + \beta + \gamma_{中文}\right)^2
			\label{eq: obj 2} \\
			\text{s.t.} \quad \boldsymbol{\alpha} \leq 9, \subeqtag \\
			\beta \geq -10, \subeqtag \\
			\gamma_{中文} \geq 8, \subeqtag[eq: constriant gamma]
		\end{align}
	\end{verbatim}
	\begin{align}
		\thinmuskip=-3mu \medmuskip=-2mu \thickmuskip=-1mu
		\max \log \left(x^2 + y^2 + z^2 + v^2 + g^2 + m^2 + k^2\right) \label{eq: obj 1} \\
		\text{s.t.} \quad x \leq 1, \subeqtag[eq: constraint x] \\
		y \leq 2, \subeqtag \\
		z \leq 4, \subeqtag
	\end{align}
	\begin{align}
		\min \left(\boldsymbol{\alpha} + \beta + \gamma_{中文}\right)^2 \label{eq: obj 2} \\
		\text{s.t.} \quad \boldsymbol{\alpha} \leq 9, \subeqtag \\
		\beta \geq -10, \subeqtag \\
		\gamma_{中文} \geq 8, \subeqtag[eq: constriant gamma]
	\end{align}
	

	$x^2 + \boldsymbol{x}^2 + y^2 + \mathrm{y}^2 + z^2 + \mathcal{B}^2 + \mathcal{M}^2 + \mathbb{I}^2 + v^2 + g^2 + m^2 + k^2$
	
	$\check{\boldsymbol{C}}^t_n + \hat{p}^j_{n, s} + \tilde{r}^t_{u,b} + \overline{i}^t_{x, y} + \acute{\alpha}^t_v + \acute{\boldsymbol{\alpha}}^t + \hat{\boldsymbol{\alpha}}^t + f^t_j + f^{"t}_j + l^t_j + l^{"t}_j + \boldsymbol{\acute{\beta}}^t + \boldsymbol{\hat{\beta}}^t$

	$\xlongequal{uvw}, \cdot, \cdots, \xLeftrightarrow{xyz}, \mathcal{N} \triangleq \left\{1, \dots, N\right\}, \times, \partial, \emptyset, \in, \subseteq, \leftarrow, \rightarrow, \leq, \geq$

	$ \alpha \beta \gamma \delta \epsilon \zeta \eta \theta \iota \kappa \lambda \mu \nu \omicron \xi \pi \rho \sigma \tau \upsilon \phi \chi \psi \omega$

	$ \boldsymbol{\alpha} \boldsymbol{\beta} \boldsymbol{\gamma} \boldsymbol{\delta} \boldsymbol{\epsilon} \boldsymbol{\zeta} \boldsymbol{\eta} \boldsymbol{\theta} \boldsymbol{\iota} \boldsymbol{\kappa} \boldsymbol{\lambda} \boldsymbol{\mu} \boldsymbol{\nu} \boldsymbol{\xi} \boldsymbol{\omicron} \boldsymbol{\pi} \boldsymbol{\rho} \boldsymbol{\sigma} \boldsymbol{\tau} \boldsymbol{\upsilon} \boldsymbol{\phi} \boldsymbol{\chi} \boldsymbol{\psi} \boldsymbol{\omega}$

	$ \Alpha \Beta \Gamma \Delta \Epsilon \Zeta \Eta \Theta \Iota \Kappa \Lambda \Mu \Nu \Xi \Omicron \Pi \Rho \Sigma \Tau \Upsilon \Phi \Chi \Psi \Omega$
	
	$ \boldsymbol{\Alpha} \boldsymbol{\Beta} \boldsymbol{\Gamma} \boldsymbol{\Delta} \boldsymbol{\Epsilon} \boldsymbol{\Zeta} \boldsymbol{\Eta} \boldsymbol{\Theta} \boldsymbol{\Iota} \boldsymbol{\Kappa} \boldsymbol{\Lambda} \boldsymbol{\Mu} \boldsymbol{\Nu} \boldsymbol{\Xi} \boldsymbol{\Omicron} \boldsymbol{\Pi} \boldsymbol{\Rho} \boldsymbol{\Sigma} \boldsymbol{\Tau} \boldsymbol{\Upsilon} \boldsymbol{\Phi} \boldsymbol{\Chi} \boldsymbol{\Psi} \boldsymbol{\Omega}$
	
	有两点需要提醒：
	\begin{itemize}
		\item \shadcmd{subeqtag[<子公式编号标签>]}的可选参数全文不可重复定义，因为它本质上还是调用的\shadcmd{label\{<编号标签>\}}。
		\item 尽管使用\shadcmd{subeqtag[<子公式编号标签>]}的可选参数指定的标签本质上是基于\shadcmd{label\{<编号标签>\}}进行的封装，TeXstudio编辑器在使用\shadcmd{ref\{<编号标签>\}}或\shadcmd{eqref\{<编号标签>\}}却不会自动弹出这些标签的选项，需要手动输入；而如果是直接用\shadcmd{label\{<编号标签>\}}指定的标签，引用时会出现在选项提示中，可以直接选择。
	\end{itemize}
	
	
	\subsection{排版及微调数学公式}

	对于不熟悉数学符号与LaTeX源码对应关系的用户，请参考\href{https://www.cnblogs.com/1024th/p/11623258.html}{\color{DarkRed}\CJKunderline*[thickness=0.5bp, format=\color{DarkRed}]{LaTeX公式手册（全网最全）@樱花赞}}，此处不再赘述。

	当某些行间公式过长以致于超出页面边界时，用户可以在\shad{equation}环境中嵌套\shad{aligned}环境将之调整为多行排版。此外，若公式超出页面边界的部分不多，也可以在\shad{equation}环境内通过调整公式中数学符号间的三种间距\shadcmd{thinmuskip}、\shadcmd{medmuskip}、\shadcmd{thickmuskip}来略微压缩公式的排版长度。这三项长度的单位只能是\shad{mu}，用户可对比公式\eqref{eq: 超长行间公式}和\eqref{eq: 微调公式中数学符号间距}的排版效果：两者的内容完全相同，但式\eqref{eq: 微调公式中数学符号间距}调整了上述三项长度。
	\begin{equation} \label{eq: 超长行间公式}
		\sum_{k=a}^{b} f(k) = \int_{a}^{b} f(x) \, dx + \frac{f(a) + f(b)}{2} + \sum_{m=1}^{\lfloor p/2 \rfloor} \frac{B_{2m}}{(2m)!} \left( f^{(2m-1)}(b) - f^{(2m-1)}(a) \right) + R_p,
	\end{equation}
	\begin{equation} \label{eq: 微调公式中数学符号间距}
		\thinmuskip=-1mu \medmuskip=-1mu \thickmuskip=0mu
		\sum_{k=a}^{b} f(k) = \int_{a}^{b} f(x) \, dx + \frac{f(a) + f(b)}{2} + \sum_{m=1}^{\lfloor p/2 \rfloor} \frac{B_{2m}}{(2m)!} \left( f^{(2m-1)}(b) - f^{(2m-1)}(a) \right) + R_p,
	\end{equation}


	\subsection{在标题中排版数学符号\texorpdfstring{$\tilde{r}^t_{u,b}, \acute{\alpha}^t_v, \check{\boldsymbol{C}}^t_n$}{示例}}

	如果需要在标题中排版数学符号，可参考本节标题的做法，使用\href{https://mirrors.tuna.tsinghua.edu.cn/CTAN/macros/latex/contrib/hyperref/doc/hyperref-doc.pdf}{\ttfamily\color{DarkRed}\CJKunderline*[thickness=0.5bp, format=\color{DarkRed}]{hyperref}}宏包（已载入）提供的\shadcmd{texorpdfstring\{<TeXstring>\}\{<PDFstring>\}}命令，该命令的具体用法参考这个帖子：\href{https://blog.csdn.net/qq_42679415/article/details/139592054}{\ttfamily\color{DarkRed}\CJKunderline*[thickness=0.5bp, format=\color{DarkRed}]{texorpdfstring使用方法}}。

	
	\subsection{排版化学方程式}

	本模板基于\href{https://mirrors.cloud.tencent.com/CTAN/macros/latex/contrib/mhchem/mhchem.pdf}{\ttfamily\color{DarkRed}\CJKunderline*[thickness=0.5bp, format=\color{DarkRed}]{mhchem}}和\href{https://mirrors.bfsu.edu.cn/CTAN/macros/generic/chemfig/chemfig-en.pdf}{\ttfamily\color{DarkRed}\CJKunderline*[thickness=0.5bp, format=\color{DarkRed}]{chemfig}}宏包来排版化学方程式、结构式和键线式等。相关命令的使用可参考宏包的官方文档或\href{https://www.luogu.com.cn/user/45443}{\color{DarkRed}\CJKunderline*[thickness=0.5bp, format=\color{DarkRed}]{@codesonic}}的文章：\href{https://www.luogu.com.cn/article/o7mlv3w8}{\color{DarkRed}\CJKunderline*[thickness=0.5bp, format=\color{DarkRed}]{用LaTeX写化学方程式}}。下方示例取自该文章，感谢。

	\ce{2H2 + O2 ->T[点燃] 2H2O}，\quad \ce{N2 + 3H2 <=>T[高温、加压][催化剂] 2NH3}
	
	\ce{SO4^2- + Ba^2+ -> BaSO4 v}，\quad \ce{2HCl + Na2CO3 -> H2O + CO2 ^ + 2NaCl}

	\ce{^{227}_{90}Th+}，\quad \ce{KCr(SO4)2 * 12H2O}，\quad \ce{C6H5-CHO}，\quad \ce{X=Y#Z}

	\ce{\chemfig{CH_3C(=[1]O)(-[7]CH_3)} + \chemfig{CN(-[2]H)} ->T[催化剂] \chemfig{CH_3(-[0]C(-[2]OH)(-[0]CN)(-[6]CH3))}}

	\ce{\chemfig{CH_3C(=[1,0.7]O)(-[7,0.7]CH_3)} + \chemfig{CN(-[2,0.7]H)} ->T[催化剂] \chemfig{CH_3(-[0,0.7]C(-[2,0.7]OH)(-[0,0.7]CN)(-[6,0.7]CH3))}}

	\chemfig{[,0.7]*6(-=(*5(-N(-H)-=(-[:30]CH_2CH_2NHCOCH_3)--))-=-(-H_3CO)=)}

	\chemfig{[,0.7]*6(-=(*5(-[0.7]N(-H)-=(-[:30,0.7](-[:330,0.7](-[:30,0.7]N(-[2,0.7]H)(-[:330,0.7](=[6,0.7]O)(-[:30,0.7])))))--))-=-(-[,0.7]O(-[1,0.7]))=)}

	
	\subsection{引用}
	
	对“公式”、“图片”、“表格”、“伪码”、“定义”、“公理”、“定理”、“命题”、“推论”、“引理”、“示例”、“假设”、“证明”等编号的引用直接用\shadcmd{ref\{<编号label>\}}即可，其中需要带括号公式编号则使用\shadcmd{eqref\{<公式label>\}}。
	
	若要对子图题编号进行完整引用直接使用\shadcmd{ref\{<子图题标签>\}}即可，\shad{nsfc}文档类默认生成形如\textcolor{DodgerBlue}{1-1(a)}的完整编号，但若指定了文档类的\shad{subfigsimple}选项，则会生成形如\textcolor{DodgerBlue}{1-1a}的完整编号；反之，若只希望单独引用子图题编号，比如在主图题结尾按编号添加子图题文本，则需要使用\shadcmd{subref\{<子图题标签>\}}，它将生成形如\textcolor{DodgerBlue}{(a)}的单独编号。
	
	对参考文献的行内引用直接使用\shadcmd{cite\{<参考文献label>\}}，以上标形式引用则使用\shadcmd{citess\{<参考文献label>\}}。
	
	参考文献的引用是基于\href{https://mirrors.zju.edu.cn/CTAN/macros/latex/contrib/natbib/natbib.pdf}{\ttfamily\color{DarkRed}\CJKunderline*[thickness=0.5bp, format=\color{DarkRed}]{natbib}}宏包实现，可单次引用多篇参考文献，届时序号会被自动排序并压缩（如果可以的话）。
	
	
	\subsection{编译参考文献}
	
	本模板提供了\textbf{“期刊论文”}、\textbf{“会议论文”}、\textbf{“专著”}、\textbf{“学位论文”}、\textbf{“报纸文章”}、\textbf{“报告”}、\textbf{“授权专利”}、\textbf{“标准”}、\textbf{“电子文献”}，共计9种文献类型的排版风格。
	
	本模板为这些文献类型定义的\shad{.bib}数据库条目\textbf{类型标识}分别为\shad{article}、\shad{inproceedings/conference}、\shad{book}、\shad{mastersthesis/phdthesis}、\shad{news}、\shad{report}、\shad{patent}、\shad{standard}、\shad{digital}。
	
	不同文档类型条目包含不同的域，以下列举了一些参考文献对应的\shad{.bib}数据，覆盖上述9种文献类型：
	
	\begin{verbatim}
	@book{教育部国家语言文字工作委员2018,
	    author={教育部国家语言文字工作委员},
	    title={通用规范汉字},
	    address={北京},
	    publisher={语文出版社},
	    year={2018},
	    language={schinese},
	}
	
	@standard{学位论文编写规范555,
	    author={全国信息与文献标准化技术委员},
	    title={学位论文编写规范},
	    number={GB/T 7713.1-2006},
	    address={北京},
	    publisher={中国标准出版社},
	    year={2007},
	    pages={17-20},
	}
	
	@article{王晓琰2019关于连续出版会议论文著录格式的探讨,
	    title={关于连续出版会议论文著录格式的探讨},
	    author={王晓琰 and 殷建芳 and 王晓峰 and 邓迎 and 杨蕾},
	    journal={学报编辑丛论},
	    number={0},
	    year={2019},
	    pages={162-165},
	    language={schinese},
	}
	
	@article{hu2014domain,
	    title={Domain decomposition method based on integral equation
	    for solution of scattering from very thin, conducting cavity},
	    author={Hu, Jun and Zhao, Ran and Tian, Mi and Zhao, Huapeng and
	    Jiang, Ming and Wei, Xiang and Nie, Zai Ping},
	    journal={IEEE Transactions on Antennas and Propagation},
	    volume={62},
	    number={10},
	    pages={5344-5348},
	    year={2014},
	    publisher={IEEE}
	}
	
	@inproceedings{bergamasco2015adopting,
	    title={Adopting an unconstrained ray model in light-field cameras
	    for 3d shape reconstruction},
	    author={Bergamasco, Filippo and Albarelli, Andrea and Cosmo, Luca
	    and Torsello, Andrea and Rodola, Emanuele and Cremers, Daniel},
	    booktitle={IEEE Conference on Computer Vision and Pattern Recognition},
	    pages={3003-3012},
	    year={2015},
	    organization={Boston, USA}
	}
	
	@article{xue2024survey,
	    title={A survey of beam management for mmWave and THz
	    communications towards 6G},
	    author={Xue, Qing and Ji, Chengwang and Ma, Shaodan and Guo, Jiajia
	    and Xu, Yongjun and Chen, Qianbin and Zhang, Wei},
	    journal={IEEE Communications Surveys \& Tutorials},
	    year={2024},
	    pages={1-41},
	    publisher={IEEE}
	}
	
	@book{罗杰斯2011,
	    author={罗杰斯},
	    title={西方文明史：问题与源头},
	    translator={潘惠霞 and 魏婧 and 杨艳 and 汤玲},
	    edition={2},
	    address={大连},
	    publisher={东北财经大学出版社},
	    year={2011},
	    pages={1-353},
	    language={schinese},
	}
	
	@book{harrington1993field,
	    title={Field computation by moment methods},
	    author={Harrington, Roger F},
	    year={1993},
	    pages={76-112},
	    edition={3},
	    address={New York},
	    publisher={Wiley-IEEE Press}
	}
	
	@digital{电子文献1,
	    author={Deverell, W and gler, D},
	    title={A companion to California history},
	    type={M/OL},
	    modifydate={2013-11-15},
	    url={http://onlinelibrary.wiley.com/doi/.ch2/summary},
	    doi={10.1002/9781444305036},
	    address={New York},
	    publisher={John Wiley \& Sons},
	    year={2013},
	    pages={21-22},
	    citedate={2014-06-24},
	}
	
	@digital{电子文献2,
	    author={Clerc, M},
	    title={Discrete particle swarm optimization: a fuzzy
	    combinatorial box},
	    type={EB/OL},
	    modifydate={2010-07-16},
	    url={http://clere.maurice.free.fr/pso/Fuzzy_Discrere_PSO/Fuzzy_DPSO.html},
	}
	
	@mastersthesis{陈念永2001毫米波细胞生物效应及抗肿瘤研究,
	    author={陈念永},
	    title={毫米波细胞生物效应及抗肿瘤研究},
	    address={成都},
	    school={电子科技大学},
	    year={2001},
	    pages={50-60},
	}
	
	@news{顾春20122,
	    author={顾春},
	    title={牢牢把握稳中求进的总基调},
	    publisher={人民日报},
	    year={2012},
	    month={03},
	    day={31},
	    number={3},
	}
	
	@report{冯西桥1997,
	    author={冯西桥},
	    title={核反应堆压力容器的{LBB}分析},
	    address={北京},
	    publisher={清华大学核能技术设计研究院},
	    year={1997},
	}
	
	@patent{肖珍新2012,
	    author={肖珍新},
	    title={一种新型排渣阀调节降温装置},
	    number={ZL201120085830.0},
	    year={2012},
	    month={04},
	    day={25},
	}

	@phdthesis{陈念永2001毫米波细胞生物效应及抗肿瘤研究无页码,
	    author={陈念永},
	    title={毫米波细胞生物效应及抗肿瘤研究（无页码测试）},
	    address={成都},
	    school={电子科技大学},
	    year={2001},
	}
	\end{verbatim}
	
	上述\shad{.bib}数据编译后的结果见本文档附上的参考文献列表。若用户漏掉了参考文献需要的强制域时，BibTeX编译会报错。在VSCode中，编译将直接中断；在TeXstudio中，编译不会中断，但log窗口会打印错误信息。遇到这类问题，应自行筛查疏漏。
	
	生成参考文献最耗费精力的是维护正确的\shad{.bib}数据库。在这之后，只需要在正文的对应位置使用以下命令即可插入完整的参考文献列表：

	\shadcmd{bibliography\{<.bst文件名>\}}
	
	\textbf{\textcolor{DarkRed}{重要提醒}：为方便演示，示例文档在上述命令之前使用了命令\shadcmd{nocite\{*\}}，该命令现处于注释状态。其作用是在参考文献列表中列出\shad{.bib}数据库中的所有参考文献条目，不论其是否在文中被引用。\textcolor{DarkRed}{正式撰写申报书时一定要确保该命令处于注释状态！！！}}
	
	\textbf{补充说明}：
	\begin{itemize}
		\item 对于某些缺少非必要信息的文献，本模板提供的\shad{.bst}文件依然可以顺利处理。比如\cite{王晓琰2019关于连续出版会议论文著录格式的探讨}这篇期刊论文缺少卷号，它仍能仅排版期号。再比如，文献\cite{电子文献2}比文献\cite{电子文献1}少了\textbf{出版地}、\textbf{出版者}等信息，依然能正常排版；但是注意，\cite{电子文献2}已经是这类文献的最简形式，不可再缺信息。
		
		\item 对中文参考文献，如果希望将它们的第四顺位及以后的作者显示为\shad{“等”}，则必须要在它们的bib条目中加入\shad{language=\{\}}域，并将值设置为\shad{schinese}。这是文献编译引擎判断该条参考文献是否是中文的唯一依据。类似的，\cite{罗杰斯2011}中的\shad{“等译”}、\shad{“2版”}均靠设置\shad{language=\{schinese\}}实现。\textbf{\textcolor{DarkRed}{虽然\shad{language}域并非强制，但对于中文文献，建议添加此域}}。
		
		\item 对电子文献，其类型众多，因此需要用户通过\shad{type=\{\}}域显式指定，如文献\cite{电子文献1}和\cite{电子文献2}；而对其他的文献类型，只要在\shad{@}符号后输入了正确的类型标识，对应的类型标签会自动生成，无需用户手动逐条添加。
	\end{itemize}
	
	%% 生成参考文献
	% \nocite{*}% 为了便于在示例文档中展示参考文献而设，正式撰写时需要注释掉
	\bibliography{ref}
	
\end{document}